\documentclass[11pt]{article}

\usepackage[margin=1in]{geometry}
\usepackage{amsmath,amssymb,amsthm,mathtools}
\usepackage{booktabs}
\usepackage{graphicx}
\usepackage{enumitem}
\usepackage{algorithm}
\usepackage{algpseudocode}
\usepackage{microtype}
\usepackage{xcolor}
\usepackage[hidelinks]{hyperref}

\newtheorem{definition}{Definition}[section]
\newtheorem{proposition}[definition]{Proposition}
\newtheorem{theorem}[definition]{Theorem}
\newtheorem{lemma}[definition]{Lemma}
\newtheorem{remark}[definition]{Remark}
\newtheorem{assumption}[definition]{Assumption}

\newcommand{\R}{\mathbb{R}}
\newcommand{\E}{\mathbb{E}}
\newcommand{\Pp}{\mathbb{P}}
\newcommand{\F}{\mathcal{F}}
\newcommand{\D}{\mathcal{D}}
\newcommand{\N}{\mathcal{N}}
\newcommand{\KL}{\operatorname{KL}}
\newcommand{\tr}{\operatorname{tr}}
\newcommand{\diag}{\operatorname{diag}}
\newcommand{\CA}{\mathrm{CA}}
\newcommand{\ind}{\mathbf{1}}

\title{Calibration-Aware Recurrent Neural Networks\\
for Multivariate Probabilistic Forecasting\\
under Distribution Shift}
\author{Jonathan Ma}
\date{Working paper\\\today}

\begin{document}

\hypersetup{pageanchor=false}
\maketitle

\begin{abstract}
Probabilistic forecasts are useful only when their stated uncertainty is credible.
Likelihood-trained recurrent neural networks may be miscalibrated, while Bayesian
model averaging is not guaranteed to remain calibrated when the likelihood, prior,
or conditional dynamics are misspecified. This paper develops the
Calibration-Aware Recurrent Neural Network (CA-RNN), adapting the
calibration-aware learning framework of Huang, Park, and
Simeone and existing differentiable regression and multivariate pre-rank
regularization to sequential multivariate forecasting. The training criterion
augments joint negative log likelihood with a differentiable discrepancy between
projected probability integral transforms (PITs) and the uniform distribution
and a temporal penalty on selected lagged PIT moments. CA-RNN is a conditional
stochastic forecasting model: its recurrent
state parameterizes a predictive distribution from which future observations can
be sampled. The primary empirical comparison is between RNN and CA-RNN, with
Bayesian variants retained as a supporting replication and VAR--GARCH as the
principal econometric benchmark. The evaluation then extends to model
misspecification and distribution shift. In a post-protocol simulation with 50
independent datasets and an isolated innovation-scale shift, CA-RNN and RNN
perform similarly at zero shift, whereas CA-RNN has lower Energy Score,
variogram score, interval score, test negative log likelihood, and analytic
projected-PIT calibration error at both shifted severities. A frozen replication on a disjoint
four-asset panel confirms a smaller projected-PIT calibration improvement, but
finds significantly worse MSE, Energy Score, interval score, and PIT temporal
dependence than the likelihood-trained RNN. A VAR with marginal GARCH$(1,1)$
scales and joint standardized-residual resampling provides substantially better
interval calibration than either neural model. The results show that directly
optimizing a calibration diagnostic can generalize to new assets without
improving the predictive distribution as judged by proper scores. Supporting
results prove that the finite calibration criterion is non-identifying, construct
dependent uniform PITs with zero covariance, and express Energy and interval-score
regret in quantities that the calibration loss does not control.
\end{abstract}

\newpage
\tableofcontents
\clearpage
\hypersetup{pageanchor=true}

\section{Introduction}

\subsection{Motivation}

A point forecast reports one possible future; a probabilistic forecast reports a
distribution of possible futures. The latter is operationally meaningful only if
events assigned probability $\alpha$ occur at approximately rate $\alpha$, while
the distribution remains sufficiently sharp to discriminate among outcomes.
Modern neural forecasters can optimize likelihood successfully without achieving
reliable finite-sample calibration. Bayesian learning represents parameter
uncertainty, but its calibration guarantees depend on idealized assumptions about
the likelihood, prior, posterior computation, and data-generating mechanism.

Huang, Park, and Simeone propose calibration-aware Bayesian neural networks
(CA-BNNs), combining a data-dependent calibration regularizer with the
data-independent KL regularizer of variational Bayesian learning
\cite{huang2024cabl}. Their experiments compare frequentist, Bayesian,
calibration-aware frequentist, and calibration-aware Bayesian classifiers; vary
the calibration weight; and report expected calibration error and reliability
diagrams. Their conclusion identifies evaluation under distribution shift as a
direction for future work. This paper transfers that methodological template to
continuous multivariate sequential prediction and makes misspecification and
distribution shift a central part of the evaluation.

\subsection{Research questions}

The primary research question is:

\begin{quote}
How does augmenting a likelihood-trained Gaussian RNN with differentiable
projected-PIT calibration penalties affect the calibration and predictive
quality of one-step-ahead multivariate financial forecasts, both across
chronological market regimes and relative to likelihood-only neural and
VAR--GARCH benchmarks?
\end{quote}

This question is decomposed into four focused empirical questions:
\begin{enumerate}[label=RQ\arabic*.]
    \item Across the prespecified calibration-weight grid, what tradeoff arises
    between projected calibration error and strictly proper predictive scores,
    marginal interval score, and point-forecast error?
    \item What does removing CA--RNN's temporal PIT-moment penalty reveal about
    its effect on held-out sequential calibration error, Energy Score, marginal
    interval score, coverage, and point-forecast error?
    \item Does CA--RNN add forecast value beyond the likelihood-only RNN and
    VAR--GARCH on the development and frozen external asset panels, and how does
    its performance relative to the RNN vary across four chronological market
    regimes?
    \item In a replicated experiment that changes only test-innovation scale,
    how do the paired CA--RNN-minus-RNN differences in analytic projected-PIT
    error, test NLL, Energy Score, variogram score, interval score, and MSE vary
    with shift severity?
\end{enumerate}

\subsection{Contributions and novelty boundary}

This study adapts differentiable projected-PIT regularization to a recurrent
conditional-density model trained on ordered multivariate observations. Its
sequential extension adds lagged PIT-moment restrictions. The empirical
contribution is a frozen comparison with likelihood-only recurrent and
VAR--GARCH forecasts, including an isolated scale-shift simulation,
chronological regimes, external-panel replication, and proper-score checks.
Neither differentiable PIT regularization for regression nor multivariate
pre-rank regularization is introduced here as a general new method; both have
direct precedents \cite{utpala2020quantile,dheur2023large,laajil2026prerank}.

\subsection{Scope and non-claims}

Finite projected calibration is weaker than complete multivariate distributional
calibration. Calibration of a finite test set is also not a guarantee of
conditional calibration at every history. We therefore avoid claims
that a small collection of projections proves joint correctness. The proposed
regularizer is an explicit, diagnostically aligned training objective whose value
must be established empirically with proper scores and out-of-sample tests.

\section{Related Work}

\subsection{Calibration-aware Bayesian learning}

The organizing reference is \cite{huang2024cabl}. It distinguishes
data-dependent calibration regularization from the data-independent prior
regularization used by variational Bayesian inference. Its CA-FNN objective adds a
differentiable approximation of expected calibration error to cross-entropy, and
its CA-BNN objective places the same calibration penalty inside the variational
expectation while retaining a KL penalty. The paper uses Gaussian variational
parameters and the reparameterization trick, and its main comparison comprises
FNN, BNN, CA-FNN, and CA-BNN. The reported experiments emphasize calibration-weight
sweeps and reliability diagrams.

\subsection{Calibration, sharpness, and proper scoring}

Probabilistic calibration concerns the statistical consistency between forecast
distributions and realized outcomes.  Gneiting, Balabdaoui, and Raftery
\cite{gneiting2007calibration} organize forecast evaluation around the principle
that a predictive distribution should be as sharp as possible subject to
calibration.  This distinction is central here: a diffuse predictive law can
achieve acceptable empirical coverage without being informative, while a sharp
but systematically incorrect law is not credible.  Calibration diagnostics must
therefore be accompanied by measures that reward the entire predictive
distribution.

Proper scoring rules provide that complement.  A scoring rule is proper when its
expected value is optimized by reporting the data-generating distribution and is
strictly proper when that optimum is unique \cite{gneiting2007scoring}.  Negative
log predictive density, the energy score, and the interval score consequently
answer a different question from a finite calibration penalty: they assess the
quality of the reported distribution rather than agreement with selected moment
conditions.  For multivariate quantities, the energy score extends the
continuous ranked probability score and supports sample-based evaluation
\cite{gneiting2008multivariate}.  It can, however, have limited sensitivity to
misspecified dependence; variogram scores were developed in part to expose such
errors \cite{scheuerer2015variogram}.  The manuscript therefore reports proper
scores alongside projected calibration, marginal intervals, and portfolio-tail
risk rather than interpreting any one diagnostic as sufficient.

More recent work formalizes conditional calibration and links forecast
evaluation to regression diagnostics \cite{gneiting2023conditional}.  This is an
important qualification for the present method.  CA--RNN enforces a finite set of
unconditional empirical restrictions on projected PITs; it does not establish
auto-calibration conditional on every forecast history.  The distinction
motivates both the non-identification results in Section~\ref{sec:theory} and the
use of genuinely out-of-sample proper-score comparisons.

\subsection{Density-forecast diagnostics under temporal dependence}

Probability integral transforms (PITs) convert a correctly specified continuous
univariate predictive distribution into uniform observations.  Diebold,
Gunther, and Tay \cite{diebold1998density} made PIT histograms and dependence
diagnostics a practical framework for evaluating financial density forecasts.
Berkowitz \cite{berkowitz2001testing} proposed likelihood-based tests after a
normal transformation of the PIT, allowing uniformity and independence to be
examined jointly.  For overlapping or multi-step forecasts, however, transformed
observations may be serially correlated even under the null.  Kn\"uppel
\cite{knuppel2015multistep} develops raw-moment tests that explicitly accommodate
this dependence.

These results explain the design of the CA--RNN penalty.  Marginal projected PIT
moments target distributional shape, while lagged PIT-product moments target one
limited form of sequential dependence.  The latter is diagnostically relevant,
but neither zero sample autocovariance nor uniformity on a finite grid proves
independence.  Consequently, sequential calibration error is treated as a
targeted diagnostic, not as a complete specification test.

\subsection{Calibration-aware neural prediction}

Modern neural networks can produce poorly calibrated probabilities even when
their predictive accuracy is high.  Guo et al.\ \cite{guo2017calibration}
document this phenomenon in classification and show the effectiveness of
post-hoc temperature scaling.  Kumar, Sarawagi, and Jain
\cite{kumar2018trainable} instead construct a differentiable kernel-based
calibration measure that can be included during training.  For continuous
outcomes, Kuleshov, Fenner, and Ermon \cite{kuleshov2018regression} propose a
post-hoc recalibration procedure for predictive cumulative distribution
functions and demonstrate it with deep regression and recurrent models.  These
papers establish that calibration can be measured or corrected in neural
systems. In-training regression calibration also has direct precedents:
quantile regularization penalizes departures from calibrated predictive
quantiles \cite{utpala2020quantile}; Dheur and Ben Taieb study differentiable
regression regularizers, including a smoothed empirical PIT-CDF discrepancy
\cite{dheur2023large}, and subsequently integrate recalibration into neural
regression training \cite{dheur2024design}. A differentiable PIT-based loss is
therefore not, by itself, a novel contribution of CA--RNN.
Related quantile-regression work explicitly optimizes a calibration--sharpness
tradeoff and provides another training-time continuous-outcome precedent
\cite{chung2021pinball}.

Huang, Park, and Simeone \cite{huang2024cabl} bring calibration directly into
Bayesian neural learning.  Their calibration-aware objective combines a
data-dependent differentiable calibration loss with a data-independent
variational KL regularizer.  CA--RNN preserves this organizing idea but changes
the statistical object: class probabilities are replaced by multivariate
one-step-ahead predictive distributions, and calibration is evaluated through
projected PIT and sequential moment restrictions.  This adaptation is not
mechanical.  Forecast observations are ordered and dependent, predictive
covariance matters, and a multivariate law cannot be validated by marginal
coverage alone.

\subsection{Multivariate pre-rank calibration and regularization}

Pre-rank functions reduce multivariate forecast--outcome pairs to scalar
diagnostics whose PIT-like distribution can be checked for uniformity. Allen,
Ziegel, and Ginsbourger develop this perspective for multivariate forecast
assessment \cite{allen2024multivariate}. More directly, recent multi-output
regression work develops differentiable pre-rank regularization for training
multivariate predictive distributions \cite{desobry2025prerank,laajil2026prerank}.
Those formulations include smoothed empirical PIT-CDF penalties; neither
projection nor smoothing is claimed here as a new general multivariate
calibration device. CA--RNN uses fixed linear projections as a narrower set of
pre-ranks. Its adaptation concerns history-dependent recurrent forecasting on
ordered data, the additional lagged PIT-moment penalty, and evaluation under
misspecification and distribution shift. Finite pooled restrictions still do
not establish calibration conditional on every forecast history.

\subsection{Probabilistic recurrent forecasting, misspecification, and shift}

DeepAR demonstrates that autoregressive recurrent networks can estimate
conditional predictive distributions by likelihood and share information across
related series \cite{salinas2020deepar}.  CA--RNN adopts the recurrent
conditional-density perspective but uses a joint multivariate Gaussian head,
then augments likelihood-based estimation
with calibration restrictions.  This separation is consequential under
misspecification.  White \cite{white1982misspecification} shows that maximum
likelihood estimation under a misspecified family generally converges to a
pseudo-true distribution rather than the data-generating law.  Calibration-aware
regularization can move the fitted distribution away from that likelihood
optimum, but such movement need not improve a strictly proper score.  The exact
regret arguments developed below formalize this tradeoff.

Distribution shift creates a second challenge.  Ovadia et al.\
\cite{ovadia2019shift} show empirically that predictive uncertainty methods that
appear adequate in distribution can deteriorate as dataset shift intensifies.
Financial returns add pronounced conditional heteroskedasticity; the GARCH model
of Bollerslev \cite{bollerslev1986garch} provides the canonical likelihood-based
mechanism for persistent time-varying variance.  These strands motivate the
manuscript's shifted-regime experiments and its VAR--GARCH benchmark.  A neural
calibration method should not receive credit merely for learning volatility that
a standard econometric model already represents directly.

\subsection{Positioning of the present study}

The literature already supplies differentiable calibration regularizers for
continuous regression and multivariate pre-rank methods, as well as recurrent
probabilistic prediction and PIT-based forecast diagnostics. The present study
asks whether adapting a projected-PIT penalty to sequential recurrent
forecasting, with an added temporal moment term, improves joint forecasts when
the model is misspecified or the data regime changes. The primary contrast is
the likelihood-trained RNN against CA--RNN; Bayesian variants support the
methodological comparison. The confirmatory design tests whether gains in the
optimized projected criterion transfer to proper scores, interval quality,
tail-risk estimates, and an external asset panel. Its mixed findings delimit
this adaptation rather than establishing a new general multivariate
calibration method.

Table~\ref{tab:related-positioning} makes the novelty boundary explicit. The
distinction is not that recurrent networks alone make calibration sequential;
it is that the loss and evaluation are applied to a filtration-indexed sequence
of conditional forecasts and add finite lagged-PIT restrictions. These
restrictions remain weaker than full conditional calibration.

\begin{table}[ht]
\centering
\caption{Positioning relative to the closest calibration literature. ``Shift''
means an explicit out-of-distribution or regime evaluation, not ordinary
held-out testing.}
\label{tab:related-positioning}
\small
\resizebox{\textwidth}{!}{%
\begin{tabular}{lcccc}
\toprule
Work & Outcome & Training penalty & Multivariate pre-rank & Sequential/shift focus \\
\midrule
Huang et al.\ \cite{huang2024cabl} & Classes & Differentiable ECE & No & No \\
Dheur--Ben Taieb \cite{dheur2023large} & Scalar regression & PIT-CDF & No & No \\
Laajil et al.\ \cite{laajil2026prerank} & Vector regression & PIT-CDF & Yes & No; i.i.d. design \\
This study & Vector time series & PIT-CDF plus lag moments & Fixed linear pre-ranks & Yes \\
\bottomrule
\end{tabular}
}
\end{table}

\section{Forecasting Setup and Calibration}

\subsection{Sequential prediction problem}

Let $(Y_t)_{t=1}^T$ be an adapted stochastic process with
$Y_t\in\R^d$ and filtration
$\F_t=\sigma(Y_1,\ldots,Y_t)$. Given a history length $H$, define
$X_t=(Y_{t-H+1},\ldots,Y_t)$. The task is to estimate a conditional predictive
distribution
\[
    P_t(\cdot)=\Pp(Y_{t+1}\in\cdot\mid\F_t)
\]
with a model $P_{\theta,t}$. Training, validation, and test blocks are ordered
chronologically. Standardization parameters are estimated on the training block
only.

\subsection{Joint Gaussian predictive head}

For the initial specification, the predictive law is
\begin{equation}
Y_{t+1}\mid\F_t,\theta
\sim \N_d\!\left(\mu_{\theta,t},\Sigma_{\theta,t}\right),
\qquad
\Sigma_{\theta,t}=L_{\theta,t}L_{\theta,t}^{\mathsf T}.
\label{eq:predictive}
\end{equation}
The network outputs the $d$ components of $\mu_{\theta,t}$ and the
$d(d+1)/2$ lower-triangular entries of $L_{\theta,t}$. An unconstrained diagonal
output $s_{t,j}$ is transformed as
\[
    L_{t,jj}=\operatorname{softplus}(s_{t,j})+\varepsilon,
    \qquad \varepsilon>0.
\]

\begin{proposition}[Validity of the predictive covariance]
For every network output, $\Sigma_{\theta,t}=L_{\theta,t}L_{\theta,t}^{\mathsf T}$
is symmetric positive definite.
\end{proposition}

\begin{proof}
Symmetry follows from transposition. Since every diagonal entry of the triangular
matrix $L_{\theta,t}$ is strictly positive, $L_{\theta,t}$ is nonsingular. Hence,
for any nonzero $x\in\R^d$,
\[
x^{\mathsf T}\Sigma_{\theta,t}x
=x^{\mathsf T}L_{\theta,t}L_{\theta,t}^{\mathsf T}x
=\|L_{\theta,t}^{\mathsf T}x\|_2^2>0.
\]
Thus $\Sigma_{\theta,t}$ is positive definite.
\end{proof}

\subsection{Negative log likelihood}

For an index set $I$, the mean joint Gaussian negative log likelihood is
\begin{align}
\mathcal L_{\mathrm{NLL}}(\theta;I)
&=-\frac{1}{|I|}\sum_{t\in I}
\log p_\theta(Y_{t+1}\mid X_t) \\
&=\frac{1}{2|I|}\sum_{t\in I}
\left[
d\log(2\pi)+2\sum_{j=1}^d\log L_{t,jj}
+e_t^{\mathsf T}\Sigma_{\theta,t}^{-1}e_t
\right],
\label{eq:nll}
\end{align}
where $e_t=Y_{t+1}-\mu_{\theta,t}$. Computation uses triangular solves rather
than an explicit matrix inverse.

\begin{proposition}[Strict propriety of the population log score]
Let $P_0(\cdot\mid\F_t)$ have density $p_0$ and let $p$ be any candidate
conditional density with common support. Then
\[
\E_{P_0}[-\log p(Y_{t+1}\mid\F_t)]
-\E_{P_0}[-\log p_0(Y_{t+1}\mid\F_t)]
=\E\!\left[\KL(P_0(\cdot\mid\F_t)\|P(\cdot\mid\F_t))\right]\geq0.
\]
Equality holds only when the conditional densities agree almost surely.
\end{proposition}

\begin{proof}
Condition on $\F_t$, subtract the two expected log scores, and recognize the
conditional KL divergence. Nonnegativity follows from Gibbs' inequality. Taking
expectations over $\F_t$ proves the result.
\end{proof}

This proposition explains why the calibration penalty must supplement rather than
replace a proper score: calibration alone can admit uninformative forecasts.

\subsection{Projected probability integral transforms}

Fix unit vectors $a_1,\ldots,a_R\in\R^d$. In the experiments they comprise the
$d$ coordinate directions, one equal-weight direction, and four random unit
directions fixed by a documented seed. Define
\begin{equation}
U_{t,r}=F_{\theta,t,r}(a_r^{\mathsf T}Y_{t+1}),
\label{eq:pit}
\end{equation}
where $F_{\theta,t,r}$ is the predictive CDF of $a_r^{\mathsf T}Y_{t+1}$.

\begin{proposition}[Projected Gaussian law]
Under \eqref{eq:predictive},
\[
a_r^{\mathsf T}Y_{t+1}\mid\F_t,\theta
\sim \N\!\left(a_r^{\mathsf T}\mu_{\theta,t},
a_r^{\mathsf T}\Sigma_{\theta,t}a_r\right),
\]
and therefore
\[
U_{t,r}=\Phi\!\left(
\frac{a_r^{\mathsf T}(Y_{t+1}-\mu_{\theta,t})}
{\sqrt{a_r^{\mathsf T}\Sigma_{\theta,t}a_r}}
\right).
\]
\end{proposition}

\begin{proof}
A linear transformation of a multivariate Gaussian is Gaussian. Its mean and
variance follow from linearity of expectation and the covariance transformation
rule. Applying its CDF gives the displayed PIT.
\end{proof}

\begin{theorem}[Sequential PIT property]
Suppose $F_{0,t,r}$ is the true, continuous conditional CDF of
$a_r^{\mathsf T}Y_{t+1}$ given $\F_t$. Then
$U_{t,r}=F_{0,t,r}(a_r^{\mathsf T}Y_{t+1})$ is uniformly distributed conditional
on $\F_t$. For a fixed $r$, the sequence $(U_{t,r})_t$ is independent and
$\operatorname{Uniform}(0,1)$.
\end{theorem}

\begin{proof}
For $u\in[0,1]$, continuity and the conditional probability integral transform
give
\[
\Pp(U_{t,r}\leq u\mid\F_t)=u.
\]
Thus $U_{t,r}$ is conditionally uniform and its conditional law does not depend on
$\F_t$. For times $t_1<\cdots<t_m$, iterated conditioning yields
\begin{align*}
\Pp(U_{t_1,r}\leq u_1,\ldots,U_{t_m,r}\leq u_m)
&=\E\left[
\ind_{\{U_{t_1,r}\leq u_1,\ldots,U_{t_{m-1},r}\leq u_{m-1}\}}
\Pp(U_{t_m,r}\leq u_m\mid\F_{t_m})
\right]\\
&=u_m\Pp(U_{t_1,r}\leq u_1,\ldots,U_{t_{m-1},r}\leq u_{m-1}).
\end{align*}
Induction gives $\prod_{j=1}^m u_j$, proving independence and uniformity.
\end{proof}

\begin{remark}
Uniformity for finitely many $a_r$ is necessary but not sufficient for equality of
multivariate distributions. If all one-dimensional projections agree, the
Cram\'er--Wold device identifies the joint distribution; a finite projection set
only provides a tractable approximation and diagnostic.
\end{remark}

Training uses the analytic Gaussian PIT above. The frozen external-panel
out-of-sample evaluation uses the common sample-rank approximation
\[
 \widehat U_{t,r}=\frac{\sum_{s=1}^{S}
 \ind\{a_r^{\mathsf T}Y_{t+1}^{(s)}<a_r^{\mathsf T}Y_{t+1}\}+1/2}{S+1},
\]
with $S=500$ predictive draws. This supports the same diagnostic for deterministic
Gaussian forecasts and Bayesian posterior-predictive mixtures, at the cost of
finite-simulation error.
The post-protocol evaluation audit additionally computes exact projected PITs
for deterministic Gaussian neural forecasts; these are reported separately
rather than substituted for the frozen sample-rank outcome.

\section{Calibration-Aware RNN Methodology}

\subsection{Recurrent architecture}

For a single-layer GRU, with elementwise product $\odot$,
\begin{align*}
r_t&=\sigma(W_{ir}Y_t+b_{ir}+W_{hr}h_{t-1}+b_{hr}),\\
z_t&=\sigma(W_{iz}Y_t+b_{iz}+W_{hz}h_{t-1}+b_{hz}),\\
n_t&=\tanh(W_{in}Y_t+b_{in}+r_t\odot(W_{hn}h_{t-1}+b_{hn})),\\
h_t&=n_t+z_t\odot(h_{t-1}-n_t).
\end{align*}
Linear heads map the final hidden state to $\mu_{\theta,t}$ and the unconstrained
entries of $L_{\theta,t}$.

\subsection{Differentiable projected calibration loss}

Let $q_g=g/(G+1)$ for $g=1,\ldots,G$. The empirical CDF of PIT values is
nondifferentiable in model parameters because it uses indicators. Replace the
indicator $\ind\{U_{t,r}\leq q_g\}$ by
\[
s_\tau(q_g-U_{t,r})
=\sigma\!\left(\frac{q_g-U_{t,r}}{\tau}\right),\qquad \tau>0.
\]
For a consecutive training batch $B$, define
\begin{equation}
\mathcal L_{\mathrm{cal}}(\theta;B)
=\frac{1}{GR}\sum_{g=1}^G\sum_{r=1}^R
\left[
\frac{1}{|B|}\sum_{t\in B}s_\tau(q_g-U_{t,r})-q_g
\right]^2.
\label{eq:cal-loss}
\end{equation}
This is a grid approximation to an integrated squared CDF discrepancy. Its
smoothed empirical PIT-CDF form belongs to the differentiable regression and
multivariate pre-rank regularization family
\cite{dheur2023large,desobry2025prerank,laajil2026prerank}; it is analogous
in purpose, but not in definition, to the classification AECE in
\cite{huang2024cabl}. The present instantiation uses fixed linear projections
and analytic Gaussian CDFs within a recurrent conditional forecast, rather
than proposing the discrepancy as a new general-purpose penalty.

\begin{proposition}[Differentiability]
If the network outputs are differentiable in $\theta$ and
$a_r^{\mathsf T}\Sigma_{\theta,t}a_r>0$, then
$\mathcal L_{\mathrm{cal}}(\theta;B)$ is differentiable in $\theta$.
\end{proposition}

\begin{proof}
The Cholesky parameterization, affine projection, positive square root, Gaussian
CDF, logistic smoother, finite sums, and squaring operations are differentiable on
the stated domain. Their composition is therefore differentiable.
\end{proof}

\begin{proposition}[Zero-temperature limit]
Fix PIT values such that $U_{t,r}\neq q_g$ for every $t,r,g$. As
$\tau\downarrow0$, \eqref{eq:cal-loss} converges to the squared grid discrepancy
formed with the empirical PIT CDF.
\end{proposition}

\begin{proof}
For nonzero $x$, $\sigma(x/\tau)\to\ind\{x>0\}$ as $\tau\downarrow0$. Every sum
is finite, so the limit passes through the averages and squares. The exclusion of
ties removes the ambiguous boundary value.
\end{proof}

\begin{proposition}[Finite-grid interpretation]
For fixed $B$, projections, grid, and $\tau>0$,
$\mathcal L_{\mathrm{cal}}(\theta;B)=0$ if and only if
\[
 \frac{1}{|B|}\sum_{t\in B}s_\tau(q_g-U_{t,r})=q_g
 \quad\text{for every }(g,r).
\]
Consequently, zero training loss establishes agreement only at the selected
smoothed CDF constraints; it does not establish uniformity between grid points,
for untested projections, conditional on the history, or out of sample.
\end{proposition}

\begin{proof}
Equation~\eqref{eq:cal-loss} is an average of nonnegative squared discrepancies.
It is zero exactly when each squared discrepancy is zero. The stated limitations
follow because the objective contains no other grid values, projections,
conditioning events, or observations outside $B$.
\end{proof}

\subsection{Temporal PIT-moment loss}

PIT uniformity does not rule out forecastable temporal structure. With
$\widetilde U_{t,r}=U_{t,r}-1/2$ and maximum lag $K$, define
\begin{equation}
\mathcal L_{\mathrm{seq}}(\theta;B)
=\frac{1}{RK}\sum_{r=1}^R\sum_{k=1}^K
\left[
\frac{1}{|B|-k}\sum_{t\in B:k\text{ available}}
\widetilde U_{t,r}\widetilde U_{t-k,r}
\right]^2.
\label{eq:seq-loss}
\end{equation}
Under the ideal conditional forecast, the population quantities targeted by
\eqref{eq:seq-loss} are zero by the sequential PIT theorem. The converse is not
true: finitely many zero autocovariances do not establish independence.

\begin{proposition}[Exact scope of the temporal PIT-moment loss]
Suppose $|B|>K$. Then $\mathcal L_{\mathrm{seq}}(\theta;B)=0$ if and only if
\[
 \frac{1}{|B|-k}\sum_{t\in B:k\text{ available}}
 (U_{t,r}-1/2)(U_{t-k,r}-1/2)=0
 \quad\text{for all }r=1,\ldots,R,\ k=1,\ldots,K.
\]
If the forecast is correctly specified so that each projected sequential PIT is
i.i.d. $\operatorname{Unif}(0,1)$, every corresponding population moment is zero.
Conversely, satisfaction of these finitely many moments does not imply temporal
independence, conditional uniformity, or correct joint dynamics.
\end{proposition}

\begin{proof}
The first statement again follows because \eqref{eq:seq-loss} is a sum of squares.
Under i.i.d. uniform PITs, $\E(U_{t,r}-1/2)=0$ and independence gives
$\E[(U_{t,r}-1/2)(U_{t-k,r}-1/2)]=0$ for $k\geq1$. The converse fails because a
finite set of second-order moment restrictions cannot characterize an entire
serial joint law; dependent processes can be uncorrelated at the tested lags.
\end{proof}

\begin{remark}
The centered product in \eqref{eq:seq-loss} uses the uniform-law mean $1/2$, not
the batch PIT mean. It is therefore a calibration moment rather than the usual
sample autocovariance unless the batch mean is exactly $1/2$.
\end{remark}

\subsection{Frequentist CA-RNN objective}

The proposed frequentist objective is
\begin{equation}
\widehat\theta_{\CA}
=\arg\min_\theta\left\{
\mathcal L_{\mathrm{NLL}}(\theta;\D)
+\lambda_{\mathrm{cal}}\mathcal L_{\mathrm{cal}}(\theta;\D)
+\lambda_{\mathrm{seq}}\mathcal L_{\mathrm{seq}}(\theta;\D)
\right\}.
\label{eq:ca-rnn}
\end{equation}
Both calibration terms in \eqref{eq:ca-rnn} define CA--RNN. The restriction
$\lambda_{\mathrm{seq}}=0$ is used only as a loss-component ablation that removes
the temporal PIT-moment penalty.

\subsection{Supporting Bayesian extension}

The Bayesian construction is retained to mirror the CA-BNN comparison and test
whether parameter averaging changes the calibration result. It is not the
primary methodological or empirical claim of this paper; the main estimand is
the frequentist RNN-versus-CA--RNN contrast.

Let every scalar neural parameter have variational posterior and prior
\[
q_\phi(\theta_j)=\N(\mu_j,\sigma_j^2),
\qquad
p(\theta_j)=\N(0,\sigma_0^2),
\qquad
\sigma_j=\operatorname{softplus}(\rho_j)+\epsilon.
\]
Sampling uses the reparameterization
\[
\theta_j=\mu_j+\sigma_j\varepsilon_j,
\qquad \varepsilon_j\sim\N(0,1).
\]
The independent-Gaussian KL term has closed form
\begin{equation}
\KL(q_\phi\|p)
=\sum_j\left[
\log\frac{\sigma_0}{\sigma_j}
+\frac{\sigma_j^2+\mu_j^2}{2\sigma_0^2}-\frac12
\right].
\label{eq:kl}
\end{equation}

The calibration-aware variational objective is
\begin{align}
\mathcal F^{\CA}(\phi)
={}&\E_{\theta\sim q_\phi}\left[
\mathcal L_{\mathrm{NLL}}(\theta;\D)
+\lambda_{\mathrm{cal}}\mathcal L_{\mathrm{cal}}(\theta;\D)
+\lambda_{\mathrm{seq}}\mathcal L_{\mathrm{seq}}(\theta;\D)
\right]\\
&+\frac{\beta}{N}\KL(q_\phi(\theta)\|p(\theta)).
\label{eq:ca-brnn}
\end{align}
This is the continuous-outcome recurrent analogue of the CA-BNN free energy in
\cite{huang2024cabl}. The factor $1/N$ matches the mean per-observation likelihood
convention. The weight $\beta$ remains a validation-selected generalized-Bayes
hyperparameter.

\begin{proposition}[Unbiased reparameterized gradient for the expectation]
Assume differentiation and expectation may be interchanged. If
$\theta^{(m)}=\mu+\sigma(\rho)\odot\varepsilon^{(m)}$ with independent
$\varepsilon^{(m)}\sim\N(0,I)$, then the Monte Carlo gradient of the expected
data-dependent part of \eqref{eq:ca-brnn} is unbiased.
\end{proposition}

\begin{proof}
Write the expectation with respect to the parameter-free distribution of
$\varepsilon$. Under the interchange assumption,
\[
\nabla_\phi\E_\varepsilon[f(\mu+\sigma(\rho)\odot\varepsilon)]
=\E_\varepsilon\nabla_\phi f(\mu+\sigma(\rho)\odot\varepsilon).
\]
The sample average of independent terms has this expectation. The KL gradient is
computed analytically from \eqref{eq:kl}.
\end{proof}

\begin{remark}[Dependent mini-batches]
Unlike independent classification observations, overlapping time-series windows
are dependent. Consecutive batches are required to evaluate
$\mathcal L_{\mathrm{seq}}$, but their nonlinear calibration loss is not asserted
to be an unbiased estimate of the full-sample calibration loss. Batch length and
history overlap are therefore methodological hyperparameters and sensitivity
targets.
\end{remark}

In the implementation, both calibration penalties are batch-local: lag pairs
crossing mini-batch boundaries are omitted. Epoch summaries average batch means,
so the final shorter batch receives the same summary weight as a full batch. The
displayed dataset-level objectives are the target criteria; stochastic
optimization uses this documented batch-local surrogate.

\subsection{Why targeted calibration need not improve proper scores}
\label{sec:theory}

This section gives a theoretical interpretation of the negative confirmatory
result. The central distinction is between satisfying a finite collection of
calibration moments and estimating the complete conditional distribution. The
former is the explicit target of CA-RNN; the latter is assessed by strictly
proper scores.

\subsubsection{Population target and finite-projection non-identification}

For a fixed forecast law $Q$ and projection $a_r$, let
$H_{Q,r}(u)=\Pp_Q(U_{t,r}\leq u)$ denote the CDF of the resulting PIT under the
true data law. The unsmoothed population counterpart of
\eqref{eq:cal-loss} on the implemented grid is
\begin{equation}
 \mathcal C_{R,G}(Q)
 =\frac{1}{RG}\sum_{r=1}^R\sum_{g=1}^G
 \{H_{Q,r}(q_g)-q_g\}^2.
 \label{eq:population-grid-calibration}
\end{equation}
It measures unconditional agreement at $RG$ points. It is not a metric on the
space of joint conditional distributions.

\begin{proposition}[Finite-grid calibration is non-identifying]
Even in one dimension with a continuous strictly increasing true CDF $F$, there
are infinitely many forecast laws $Q\neq P$ for which
$\mathcal C_{1,G}(Q)=0$ on any fixed finite grid.
\end{proposition}

\begin{proof}
Choose any strictly increasing bijection $h:[0,1]\to[0,1]$ that fixes every grid
point, $h(q_g)=q_g$, but is not the identity between at least two adjacent grid
points. Define forecast law $Q$ by its CDF $G(y)=h(F(y))$. If $Y\sim F$, then
$G(Y)=h(F(Y))$ and, because $F(Y)$ is uniform,
\[
 \Pp\{G(Y)\leq q_g\}=h^{-1}(q_g)=q_g.
\]
Thus every grid discrepancy is zero, while $G\neq F$. Infinitely many such
monotone deformations $h$ can be constructed between grid points.
\end{proof}

The multivariate case has two additional sources of non-identification: only a
finite set of directions is inspected, and the moments are averaged across those
directions. The Cram\'er--Wold device would identify a joint law if all
one-dimensional projected distributions agreed, but it does not turn a finite
projection set into an identifying criterion.

\begin{proposition}[Consistency of fixed calibration moments]
Fix $Q$, $R$, $G$, and $\tau>0$. If the PIT process is stationary and ergodic,
then the empirical smoothed grid moments converge almost surely to their
population counterparts. The empirical lagged moments in
\eqref{eq:seq-loss} likewise converge almost surely for every fixed tested lag,
provided their products are integrable.
\end{proposition}

\begin{proof}
The smoothed indicators are bounded measurable functions of the PIT process, so
Birkhoff's ergodic theorem applies to each sample average. The lagged products
are measurable functions of the finite-dimensional shifted process and are
integrable by assumption. There are finitely many projections, grid points, and
lags, so convergence holds jointly for all implemented moments.
\end{proof}

This result concerns a fixed forecast rule under one stationary law. It does not
provide uniform convergence over neural parameters, correct batch-local
optimization bias, or transfer the moments to a shifted test law. Those missing
implications are precisely relevant to the observed training-to-test gap.

\subsubsection{Zero PIT autocovariance does not imply independence}

The limitation of the sequential penalty is constructive, not merely technical.

\begin{proposition}[Dependent uniform PITs with zero covariance]
There exist dependent $U,V\sim\operatorname{Unif}(0,1)$ satisfying
$\E[(U-1/2)(V-1/2)]=0$.
\end{proposition}

\begin{proof}
Let $f(u)=6u^2-6u+1$ and define a bivariate density on $[0,1]^2$ by
\[
 c_\varepsilon(u,v)=1+\varepsilon f(u)f(v),
 \qquad 0<|\varepsilon|<1.
\]
Because $\int_0^1 f(u)\,du=0$, both marginals are uniform. The chosen range of
$\varepsilon$ keeps the density nonnegative. Since $c_\varepsilon$ is not
constant, $U$ and $V$ are dependent. However, symmetry of $f$ about $1/2$ gives
\[
 \int_0^1(u-1/2)f(u)\,du=0,
\]
and hence
\[
 \E[(U-1/2)(V-1/2)]
 =\varepsilon\left\{\int_0^1(u-1/2)f(u)\,du\right\}^2=0.
\]
\end{proof}

Consequently, even population minimization of every moment represented by a
finite-lag version of \eqref{eq:seq-loss} cannot establish PIT independence.
The confirmatory increase in PIT autocorrelation is stronger still: it shows that
the estimated training moments did not generalize as lower corresponding test
moments.

\subsubsection{Energy Score regret is not controlled by calibration loss}

For distributions with finite first moments, define
\[
 \operatorname{ES}(Q,y)
 =\E\|X-y\|_2-\frac12\E\|X-X'\|_2,
 \qquad X,X'\stackrel{\mathrm{iid}}\sim Q.
\]

\begin{proposition}[Energy Score regret]
If $Y\sim P$, then
\begin{equation}
 \E_P\operatorname{ES}(Q,Y)-\E_P\operatorname{ES}(P,Y)
 =\frac12\mathcal E(P,Q)\geq0,
 \label{eq:energy-regret}
\end{equation}
where
\[
 \mathcal E(P,Q)=2\E\|X-Y\|_2-\E\|X-X'\|_2-\E\|Y-Y'\|_2
\]
is the energy distance. Equality holds only when $P=Q$.
\end{proposition}

\begin{proof}
Expand the two expected Energy Scores using independent
$X,X'\sim Q$ and $Y,Y'\sim P$, then subtract. The remaining expression is one
half of the energy distance. Its nonnegativity and identity of indiscernibles are
the defining properties of energy distance on distributions with finite first
moments.
\end{proof}

Equation~\eqref{eq:energy-regret} shows that Energy Score evaluates the complete
joint distribution, whereas \eqref{eq:population-grid-calibration} evaluates
only finitely many scalar constraints. No inequality in the CA-RNN objective
bounds $\mathcal E(P,Q)$ by $\mathcal C_{R,G}(Q)$.

\begin{proposition}[Calibration improvement need not improve a strictly proper score]
For any fixed finite PIT grid, there exist a true continuous distribution $P$ and
forecasts $Q_0,Q_1$ such that
\[
 \mathcal C_{1,G}(Q_1)<\mathcal C_{1,G}(Q_0),
 \qquad
 \E_P S(Q_1,Y)>\E_P S(Q_0,Y)
\]
for any strictly proper score $S$ that is finite and continuous along sufficiently
small CDF perturbations of $P$.
\end{proposition}

\begin{proof}
Use the preceding monotone construction to choose $Q_1\neq P$ with zero
finite-grid calibration loss. Strict propriety gives a positive score regret for
$Q_1$. Now choose $Q_0\neq P$ as an arbitrarily small monotone perturbation of
$P$ that does not fix at least one grid point. Then
$\mathcal C_{1,G}(Q_0)>0$, while continuity makes its proper-score regret
arbitrarily close to zero. Choose the perturbation small enough that its regret
is below that of $Q_1$. The two displayed inequalities follow.
\end{proof}

This proposition formalizes the main empirical result: a lower projected-PIT
error is not an ordering of forecast distributions and cannot, by itself, imply a
lower Energy Score.

\subsubsection{Why interval score can also worsen}

Let $F$ be the true marginal CDF, $\alpha$ the nominal miscoverage probability,
$\tau=\alpha/2$, and $(\ell^*,u^*)=(F^{-1}(\tau),F^{-1}(1-\tau))$ the optimal
central interval. For a candidate interval $(\ell,u)$, the expected interval-score
regret has an exact quantile representation.

\begin{proposition}[Interval-score regret]
For continuous $F$ with finite first moment,
\begin{align}
 &\E_F\operatorname{IS}_\alpha(\ell,u;Y)
 -\E_F\operatorname{IS}_\alpha(\ell^*,u^*;Y) \\
 &\quad=\frac{1}{\tau}
 \left\{
 \int_{\ell^*}^{\ell}[F(x)-\tau]\,dx
 +\int_{u^*}^{u}[F(x)-(1-\tau)]\,dx
 \right\}\geq0.
 \label{eq:interval-regret}
\end{align}
\end{proposition}

\begin{proof}
The interval score equals $\tau^{-1}$ times the sum of pinball losses at levels
$\tau$ and $1-\tau$. The derivative of expected pinball loss with respect to a
candidate quantile $q$ is $F(q)-p$ at level $p$. Integrating this derivative from
each true quantile to its candidate value yields \eqref{eq:interval-regret}.
Monotonicity of $F$ makes both oriented integrals nonnegative.
\end{proof}

The CA-RNN loss averages 19 grid levels and multiple directions and is optimized
with smoothing on dependent training batches. A reduction in that average need
not reduce coordinate-specific errors at the $0.05$ and $0.95$ quantiles on a
shifted test distribution. Equation~\eqref{eq:interval-regret} shows that any such
tail-quantile displacement creates positive expected interval-score regret. This
explains how projected calibration can improve while the interval score worsens:
the scalar calibration average falls, but errors at the particular marginal tail
functionals priced by the interval score increase or the added width is not offset
by sufficiently smaller miss penalties.

\subsection{Training algorithms}

\begin{algorithm}[ht]
\caption{Frequentist CA-RNN training}
\label{alg:ca-rnn}
\begin{algorithmic}[1]
\Require Ordered observations $\D$; history $H$; projections $A$; weights
$\lambda_{\mathrm{cal}},\lambda_{\mathrm{seq}}$; epoch limit $E$;
gradient threshold $c$
\Ensure Fitted parameters $\widehat\theta$ selected without test information
\State Split $\D$ chronologically into $\D_{\mathrm{tr}},\D_{\mathrm{val}},
\D_{\mathrm{te}}$
\State Fit standardization on $\D_{\mathrm{tr}}$ and construct length-$H$ windows
\State Initialize $\theta$ and fix $A=(a_1,\ldots,a_R)^{\mathsf T}$
\For{$e=1,\ldots,E$}
    \For{consecutive mini-batch $B\subset\D_{\mathrm{tr}}$}
        \State $(\mu_{\theta,t},L_{\theta,t})_{t\in B}
        \gets \Call{CARNNForward}{B,\theta}$
        \State Compute $\mathcal L_{\mathrm{NLL}}$ by \eqref{eq:nll} and
        projected PITs by \eqref{eq:pit}
        \State Compute $\mathcal L_{\mathrm{cal}}$ by \eqref{eq:cal-loss} and
        $\mathcal L_{\mathrm{seq}}$ by \eqref{eq:seq-loss}
        \State $\mathcal J\gets\mathcal L_{\mathrm{NLL}}
        +\lambda_{\mathrm{cal}}\mathcal L_{\mathrm{cal}}
        +\lambda_{\mathrm{seq}}\mathcal L_{\mathrm{seq}}$
        \State $g\gets\Call{ClipGradNorm}{\nabla_\theta\mathcal J,c}$
        \State Update $\theta$ with AdamW using $g$
    \EndFor
    \State Evaluate the prespecified selection score on $\D_{\mathrm{val}}$
    \State Save $\theta$ if the validation score improves
    \If{the patience criterion is met} \State \textbf{break} \EndIf
\EndFor
\State Restore the selected checkpoint $\widehat\theta$
\State Evaluate $\widehat\theta$ once on $\D_{\mathrm{te}}$
\end{algorithmic}
\end{algorithm}

For fixed-budget causal loss ablations, checkpoint selection and the patience test
in Algorithm~\ref{alg:ca-rnn} are disabled, and all specifications are compared
after the same number of epochs.

\begin{algorithm}[ht]
\caption{Variational calibration-aware Bayesian RNN training}
\label{alg:ca-brnn}
\begin{algorithmic}[1]
\Require Ordered observations $\D$; projections $A$; prior scale $\sigma_0$;
Monte Carlo count $M$; loss weights $\lambda_{\mathrm{cal}},
\lambda_{\mathrm{seq}},\beta$; epoch limit $E$
\Ensure Variational parameters $\widehat\phi=(\widehat\mu,\widehat\rho)$
\State Construct chronological splits and training-only standardization as in
Algorithm~\ref{alg:ca-rnn}
\State Initialize $\phi=(\mu,\rho)$ and set
$\sigma(\rho)=\operatorname{softplus}(\rho)+\epsilon$
\For{$e=1,\ldots,E$}
    \For{consecutive mini-batch $B\subset\D_{\mathrm{tr}}$}
        \For{$m=1,\ldots,M$}
            \State Draw $\varepsilon^{(m)}\sim\N(0,I)$ and set
            $\theta^{(m)}\gets\mu+\sigma(\rho)\odot\varepsilon^{(m)}$
            \State Use $\theta^{(m)}$ consistently at every recurrent time step
            \State Compute $\mathcal J^{(m)}_{\mathrm{data}}\gets
            \mathcal L_{\mathrm{NLL}}^{(m)}+\lambda_{\mathrm{cal}}
            \mathcal L_{\mathrm{cal}}^{(m)}+\lambda_{\mathrm{seq}}
            \mathcal L_{\mathrm{seq}}^{(m)}$
        \EndFor
        \State $\widehat{\mathcal F}^{\CA}\gets M^{-1}\sum_{m=1}^M
        \mathcal J^{(m)}_{\mathrm{data}}+(\beta/N)\KL(q_\phi\Vert p)$
        \State Differentiate $\widehat{\mathcal F}^{\CA}$, clip its gradient,
        and update $\phi$ with AdamW
    \EndFor
    \State Evaluate validation NLL at posterior-mean weights and update the
    selected checkpoint
    \If{the patience criterion is met} \State \textbf{break} \EndIf
\EndFor
\State Restore $\widehat\phi$
\For{each test origin $t$ and predictive draw $s$}
    \State Draw $\theta^{(s)}\sim q_{\widehat\phi}$ and then
    $Y_{t+1}^{(s)}\sim p_{\theta^{(s)}}(\cdot\mid\F_t)$
\EndFor
\State Evaluate the posterior predictive mixture on $\D_{\mathrm{te}}$
\end{algorithmic}
\end{algorithm}

Algorithm~\ref{alg:ca-brnn} is the recurrent continuous-outcome adaptation of the
CA-BNN procedure in \cite{huang2024cabl}.

\section{Model Misspecification and Distribution Shift}

\subsection{Distinguishing misspecification from shift}

Model misspecification occurs when the true conditional law $P_0(Y_{t+1}\mid
\F_t)$ is outside the fitted family, even if the data law is stationary.
Distribution shift occurs when the evaluation law differs from the development
law. The two can coexist: a Gaussian predictive head is misspecified under
heavy-tailed innovations, while a transition from Gaussian training innovations
to heavy-tailed test innovations is also a distribution shift.

Let $P_0$ denote the reference law and let $P_{\delta,s}$ denote stress family $s$
at severity $\delta\geq0$, with $P_{0,s}=P_0$. For metric $M$, define degradation
\begin{equation}
\Delta_M^{(m,s)}(\delta)
=M(P_{\delta,s},\widehat P_m)-M(P_{0,s},\widehat P_m),
\label{eq:degradation}
\end{equation}
where $m$ indexes the forecasting method. For metrics in which smaller is better,
positive $\Delta_M$ denotes deterioration. Coverage is reported as absolute error
from the nominal level before differencing.

\subsection{Stress families}

The following stress families distinguish mechanisms; only the first has been
executed as a replicated isolated-factor experiment, and none has been executed
as a full factorial study:
\begin{enumerate}
    \item \textbf{Scale shift:} innovation covariance is multiplied by a severity
    function while dependence is held fixed.
    \item \textbf{Dependence shift:} correlations or a dynamic correlation state
    change while marginal variances are controlled.
    \item \textbf{Tail shift:} Gaussian innovations transition to Student-$t$,
    mixtures, or asymmetric downside contamination.
    \item \textbf{Volatility misspecification:} ARCH/GARCH or stochastic-volatility
    dynamics violate the Gaussian head's conditionally homoskedastic approximation.
    \item \textbf{Mean-dynamics shift:} nonlinear interactions or an abrupt change
    in recurrent coefficients alter the conditional location.
\end{enumerate}

The initial developmental curve uses common random numbers and jointly changes
scale and common correlation. Models are selected without examining its stressed
test metrics and are frozen after fitting at $\delta=0$. The completed
post-protocol experiment isolates scale: 50 independent reference datasets are
generated, paired RNN and CA--RNN models are fitted on each unshifted training
history, and the same fitted models are evaluated at scale severities
$\delta\in\{0,0.5,1\}$. The innovation standard deviation becomes $1+\delta$
times its reference value only in the test partition; correlation, transition
coefficients, and innovation family remain fixed. Tail, volatility, and
mean-dynamics regimes appear in separate developmental simulations, not in this
isolated-factor experiment.

\subsection{Robustness estimands}

For paired realization $b$, method $m$, comparator $c$, and metric $M$, define
\[
D_b^{m,c}(\delta)=M_b^m(\delta)-M_b^c(\delta).
\]
For the completed isolated-scale experiment with $B=50$ independent DGP
realizations, the summary is
\[
\overline D^{m,c}(\delta)=\frac1B\sum_{b=1}^B D_b^{m,c}(\delta),
\qquad
\operatorname{MCSE}(\overline D)=
\frac{s_D(\delta)}{\sqrt B}.
\]
One distinct neural seed is assigned to each DGP realization; thus the reported
MCSE includes both DGP and optimization variation but does not identify their
separate contributions. Failure frequencies and runtimes are reported rather
than discarded implementation details. The comparison is developmental because
the model weights and shift family were selected after earlier exploratory
analyses; it is not an additional frozen confirmatory replication.

\section{Experimental Design}

\subsection{Primary comparison and supporting Bayesian replication}

The design follows the comparison logic of \cite{huang2024cabl}, but the first
and third rows define the primary frequentist comparison. The Bayesian rows are
supporting replications:
\begin{table}[ht]
\centering
\caption{Mapping from calibration-aware classification to forecasting.}
\begin{tabular}{lll}
\toprule
Source-paper model & Forecasting model & Objective components \\
\midrule
FNN & RNN & NLL \\
BNN & BRNN & expected NLL $+$ KL \\
CA-FNN & CA-RNN & NLL $+$ projected-PIT calibration \\
CA-BNN & CA-BRNN & expected NLL/calibration $+$ KL \\
\bottomrule
\end{tabular}
\end{table}

Architectures, histories, splits, optimizers, stopping rules, and predictive heads
are matched wherever the Bayesian formulation permits. The same projection matrix
is used for all models in a comparison.

\subsection{Conventional benchmark ladder}

The neural comparisons are supplemented by three increasingly dynamic
probabilistic benchmarks. A centered historical bootstrap samples joint training
returns and contains neither a dynamic conditional mean nor time-varying scale. A
Gaussian VAR estimates a linear conditional mean and a fixed innovation
covariance. The strongest conventional benchmark combines that VAR mean with
separate marginal GARCH$(1,1)$ filters,
\[
 h_{j,t}=\omega_j+\alpha_j e_{j,t-1}^2+\beta_j h_{j,t-1},
 \qquad \omega_j>0,\quad \alpha_j,\beta_j\geq0,\quad
 \alpha_j+\beta_j<1,
\]
and samples entire standardized-residual vectors
$z_t=(e_{1,t}/\sqrt{h_{1,t}},\ldots,e_{d,t}/\sqrt{h_{d,t}})$ from the training
period. Its one-step samples are
\[
 Y_t^{(s)}=\widehat\mu_t+\operatorname{diag}(\sqrt{\widehat h_t})z_{I_s},
\]
where $I_s$ is a common row index across margins. This preserves empirical
contemporaneous dependence and tail shape while allowing marginal conditional
volatility to evolve. The VAR lag is selected from $\{1,5,10,20\}$ by validation
one-step MSE; parameters and the residual resampling distribution use training
data only.

Model checking for the selected benchmark is also training-only. VAR dynamic
stability is assessed by the spectral radius of its companion matrix. GARCH
admissibility is checked through optimizer convergence, nonnegative parameters,
$\alpha_j+\beta_j<1$, and the frequency with which filtered variances reach the
numerical floor. Ljung--Box statistics at lags 5, 10, and 20 are computed for
standardized residuals and their squares to diagnose remaining conditional-mean
and volatility dependence. Standardized-residual skewness, kurtosis,
Jarque--Bera tests, and the residual correlation matrix assess the rationale for
resampling empirical residual vectors jointly. The Ljung--Box and Jarque--Bera
$p$-values are descriptive because they do not account for VAR--GARCH parameter
estimation or multiplicity.

\subsection{Frozen external-panel replication}

The confirmatory protocol is stored separately in
\texttt{docs/confirmatory\_protocol.md} before outcome inspection. It fixes a
disjoint MSFT--BAC--CVX--COST panel, the date boundary, transformations, split,
architecture, loss weights, training and forecast seeds, primary outcomes,
baseline definitions, bootstrap procedure, and failure reporting. The sole
primary contrast is CA-RNN minus RNN across ten matched optimization
seeds; conventional-baseline comparisons are secondary. Any post-outcome change
creates an exploratory analysis rather than replacing the frozen result.
The archived protocol retains the model label used when it was frozen; that label
refers to the identical projected-plus-temporal objective now called CA--RNN.

\subsection{Calibration-weight experiment}

Analogously to the source paper's plot of ECE and accuracy against $\lambda$, this
study plots projected-PIT calibration error and predictive performance against
$\lambda_{\mathrm{cal}}$. The paired panels will contain:
\begin{enumerate}
    \item projected calibration discrepancy and reliability curves;
    \item Energy Score, NLL where comparable, RMSE, interval score, and width.
\end{enumerate}
The result of interest is a Pareto relationship, not calibration alone. Separate
curves are required for CA-RNN and CA-BRNN.

\subsection{Temporal-penalty ablation}

The minimum ablation is:
\begin{table}[ht]
\centering
\caption{Loss-component ablation.}
\begin{tabular}{lccc}
\toprule
Specification & NLL & $\mathcal L_{\mathrm{cal}}$ & $\mathcal L_{\mathrm{seq}}$ \\
\midrule
RNN & \checkmark & & \\
CA-RNN (temporal penalty removed) & \checkmark & \checkmark & \\
CA-RNN & \checkmark & \checkmark & \checkmark \\
\bottomrule
\end{tabular}
\end{table}
The Bayesian versions repeat this comparison with the KL term always present.
Additional ablations vary coordinate-only versus augmented projections, smoothing
temperature, batch length, lag order, prior scale, and $\beta$.

To prevent checkpoint selection from masquerading as a calibration effect, the
primary loss-component ablation uses identical initial parameters, chronological
mini-batches, optimization budgets, and predictive random numbers, with early
stopping disabled. A secondary operational comparison may use early stopping, but
all methods must then be selected by the same validation NLL rule.

\subsection{Controlled simulations}

Development simulations include Gaussian, heavy-tailed, heteroskedastic, and
nonlinear regimes. A future confirmatory experiment should use independently
generated data for every regime, sample size, shift severity, and DGP seed, with
the following fixed in advance:
\begin{itemize}
    \item sample sizes and number of independent realizations;
    \item neural-seed subset and number of optimization repetitions;
    \item calibration and KL grids;
    \item projection directions and smoothing grid;
    \item primary metrics and failure definitions; and
    \item rules for early stopping and invalid numerical runs.
\end{itemize}

\subsection{Financial application}

The application uses adjusted daily prices for AAPL, JPM, XOM, and WMT beginning
in 2007 and converts them to aligned log returns. The four assets provide a small,
computationally manageable multivariate system with heterogeneous sector exposure.
All splits are chronological. Development selected
$\lambda_{\mathrm{cal}}=100$ and $\lambda_{\mathrm{seq}}=2000$ before the
multi-seed and rolling-regime analyses. The focused frequentist comparison uses a
32-unit GRU, a 20-day history, a joint Gaussian head, 500 forecast draws, and
validation-NLL early stopping.

Standardization is fitted on each training partition. Validation contexts carry
the final 20 training observations, and test contexts carry the final 20
validation observations; every evaluation target therefore retains an available
historical context without using future information. RMSE, Energy Score, interval
width, and interval score are reported in training-standardized units. Coverage
and PIT diagnostics are invariant to the componentwise affine transformation.
The two penalty weights were chosen after inspecting developmental simulation
and ablation outcomes, rather than by a pristine validation-only or preregistered
selection procedure.

The rolling analysis uses expanding estimation samples and four non-overlapping
test periods: 2017--2019 (pre-pandemic), 2020--2021 (pandemic), 2022--2023
(inflation and monetary tightening), and 2024--2025 (recent). Each window has a
strictly subsequent validation period and five matched optimization seeds. These
labels describe calendar periods rather than uniquely identified causal regimes.
Because the full panel was observed during development and the final two windows
overlap the initial test period, this is temporal robustness evidence rather than
a pristine preregistered confirmatory experiment.

\subsection{Status of the evidence}

The lambda sweep, ablation, original-panel ten-seed study, rolling windows,
single-realization mixed-shift curve, and replicated isolated-scale experiment
are developmental or internal-validation evidence. They
condition on a specification informed by earlier results. The external-panel
replication is confirmatory for its single primary contrast: CA-RNN
minus RNN on MSFT, BAC, CVX, and COST under the protocol frozen in
\texttt{docs/confirmatory\_protocol.md}. Its traditional-baseline comparisons
were prespecified but secondary. Inferential intervals quantify test-origin and
optimization-seed uncertainty under the implemented design; they do not include
uncertainty over asset populations, architectures, or alternative research
decisions.

\subsection{Paired uncertainty design}

For model $m$, comparator $c$, seed $s$, and forecast origin $t$, let
$d_{s,t}^{m,c}$ be a paired additive-score difference. The estimand is
\[
 \widehat\Delta^{m,c}=\frac{1}{S}\sum_{s=1}^S\frac{1}{T}
 \sum_{t=1}^T d_{s,t}^{m,c}.
\]
The hierarchical bootstrap resamples optimization seeds and, within each selected
seed, matched forecast origins in moving blocks of 20 trading days. Projected-PIT
calibration, coverage error, and PIT dependence are recomputed in every bootstrap
draw. The initial financial analysis uses 2,000 draws and each rolling regime uses
1,000. Raw coverage and width remain descriptive; the inferential coverage loss is
\[
 E_{\mathrm{cov}}=|\widehat C_{0.90}-0.90|.
\]
Negative candidate-minus-comparator differences favor the candidate for loss
metrics. An interval containing zero is called unresolved, not evidence of
equality. Inference conditions on the architecture, tuning path, projection set,
assets, and observed calendar windows.

\subsection{Metrics and figures}

The tables report mean performance, paired differences, and uncertainty
intervals. The diagnostic figures include:
\begin{enumerate}
    \item calibration and predictive metrics versus $\lambda_{\mathrm{cal}}$;
    \item nominal-versus-empirical coverage curves;
    \item projected PIT histograms or CDF reliability curves;
    \item PIT autocorrelation diagnostics;
    \item degradation curves versus shift severity; and
    \item a calibration--sharpness or calibration--Energy Score frontier.
\end{enumerate}

\section{Results}

\subsection{Implementation validation}

The implementation tests cover tensor dimensions, covariance
positive-definiteness, loss evaluation, forecasting, chronological data handling,
traditional baselines, and paired-bootstrap identities. Thirty-one automated
tests pass. The notebooks also
completed the simulation, shift, financial, optimization-seed, and rolling-regime
workflows. Successful execution establishes implementation consistency but is not
itself statistical validation.

\subsection{Primary comparison and supporting Bayesian results}

The refreshed matched nonlinear simulation showed that calibration-aware
objectives alter the forecast distribution but do not yield uniform dominance.
Relative to RNN, CA-RNN slightly reduced RMSE (0.9368 to 0.9334), Energy Score
(1.2354 to 1.2348), and mean absolute projected-PIT autocorrelation (0.0733 to
0.0674). It increased interval score (3.8771 to 3.9297) and projected-PIT
calibration error (0.001224 to 0.001455), despite improving nominal coverage
from 0.9167 to 0.9375 only in the sense of producing wider intervals; both
coverage rates exceed the 0.90 target. Bayesian averaging did not remove this
calibration--score tradeoff: CA-BRNN had worse RMSE, Energy Score, interval
score, and projected calibration than BRNN, although its projected-PIT temporal
dependence was slightly lower. RNN and CA-RNN therefore form the focused
empirical comparison; BRNN and CA-BRNN remain supporting replications of the
source paper's model matrix. These are single-realization developmental results
and are not treated as inferential evidence.

\subsection{Calibration-weight sensitivity}

The exploratory sweep over
$\lambda_{\mathrm{cal}}\in\{0,1,10,50,100,250,500\}$ displayed a Pareto
tradeoff: projected calibration improved through the intermediate penalty range,
whereas likelihood and proper-score performance deteriorated. The value 100 was
retained as a developmental compromise rather than selected to maximize a final
test result. Calibration improvement alone is therefore not interpreted as a
better probabilistic forecast.

\subsection{Temporal-penalty ablation}

The temporal term has regime- and design-dependent effects. In the refreshed
single-realization heteroskedastic simulation, removing the temporal penalty
gave Energy Score 1.2101, interval score 3.9236, RMSE 0.9222, projected-PIT
calibration error 0.000770, and mean absolute projected-PIT autocorrelation
0.05182. The default CA-RNN temporal weight of 2,000 instead gave 1.2153,
3.9409, 0.9258, 0.001046, and 0.05132, respectively. Thus the temporal term
slightly reduced the diagnostic it directly targets, but worsened the proper
scores, point error, projected uniformity, and absolute nominal-coverage error
in this controlled run. Increasing the weight further continued to reduce PIT
autocorrelation while generally worsening the other metrics.

The ten-seed financial ablation has a different ordering. There, CA-RNN improved
MSE, Energy Score, interval score, nominal-coverage error, and projected-PIT
calibration relative to temporal-penalty removal for every matched seed. It
widened intervals, but simultaneous coverage and interval-score improvements
show that the result is not merely unpenalized widening. Improvements in PIT
temporal dependence were smaller and generally unresolved. Taken together, the
experiments show that the penalty changes serial calibration behavior and can
repair costs induced by marginal calibration in financial data, but does not
provide a general predictive improvement or establish PIT independence.

\subsection{Misspecification and distribution shift}

The replicated scale-only experiment completed all 50 independent DGP
realizations and 100 neural fits without recorded failures. Table~\ref{tab:isolated-scale}
reports paired CA--RNN-minus-RNN differences, with Monte Carlo standard errors
in parentheses. At zero shift, the Energy Score difference is essentially zero:
$-0.00011$ with MCSE $0.00079$. Its interval-score difference is $+0.00311$
(MCSE $0.00435$), and analytic projected-PIT calibration differs by
$-0.000025$ (MCSE $0.000024$). This is nominal near-parity, not evidence of
uniform CA--RNN dominance: its zero-shift test NLL is higher by $0.00629$
(MCSE $0.00277$) and its variogram score is higher by $0.000414$
(MCSE $0.000107$).

At severity $0.5$, CA--RNN has lower Energy Score by $0.01624$ (MCSE $0.00152$),
interval score by $0.29711$ (MCSE $0.01749$), variogram score by $0.00764$
(MCSE $0.00040$), analytic projected-PIT calibration error by $0.00124$
(MCSE $0.000065$), and test NLL by $0.22506$ (MCSE $0.01534$). At severity
$1$, the corresponding advantages increase to $0.02924$, $0.55844$, $0.01543$,
$0.00178$, and $0.57664$, respectively. Component-averaged MSE also favors
CA--RNN under both shifts, but by smaller amounts: $-0.00644$ (MCSE $0.00260$)
and $-0.01485$ (MCSE $0.00429$). All absolute proper scores and calibration
errors worsen as scale rises for both methods; the finding is less degradation
for CA--RNN, not immunity to shift. The signs of the principal shifted effects
persist in the 10-, 25-, and 50-realization prefixes, although the earliest
prefix gives somewhat larger effect magnitudes. No inference is made about
correlation-only, tail, or mean-dynamics shifts from this experiment.

\begin{table}[ht]
\centering
\caption{Isolated Gaussian innovation-scale shift over 50 independent paired
realizations. Each entry is the mean CA--RNN-minus-RNN difference with its realization-level
Monte Carlo standard error in parentheses; negative values favor CA--RNN.}
\label{tab:isolated-scale}
\small
\begin{tabular}{lrrrr}
\toprule
Severity & Energy Score & Variogram & Interval score & Analytic PIT error \\
\midrule
$0$ & $-0.00011\,(0.00079)$ & $+0.000414\,(0.000107)$ & $+0.00311\,(0.00435)$ & $-0.000025\,(0.000024)$ \\
$0.5$ & $-0.01624\,(0.00152)$ & $-0.007638\,(0.000400)$ & $-0.29711\,(0.01749)$ & $-0.001241\,(0.000065)$ \\
$1$ & $-0.02924\,(0.00215)$ & $-0.015428\,(0.000766)$ & $-0.55844\,(0.02850)$ & $-0.001781\,(0.000090)$ \\
\bottomrule
\end{tabular}
\end{table}

\begin{figure}[ht]
\centering
\IfFileExists{../results/isolated_scale_shift/paired_shift_profiles.png}{%
\includegraphics[width=0.84\textwidth]{../results/isolated_scale_shift/paired_shift_profiles.png}%
}{\fbox{Run notebook 14 to generate the isolated-scale paired profiles.}}
\caption{Paired CA--RNN-minus-RNN differences across scale severity. Error bars
show $\pm1.96$ Monte Carlo standard errors across independent DGP realizations;
they do not cover uncertainty over alternative DGP families or tuning choices.}
\label{fig:isolated-scale}
\end{figure}

The frozen-model experiment jointly increased innovation scale and common
correlation from severity zero to one. All specifications deteriorated. At zero,
RNN had Energy Score 1.2022, interval score 3.8595, and projected calibration
error 0.000866, compared with 1.2039, 3.9311, and 0.001345 for CA-RNN.
At maximum severity, the ordering reversed: RNN obtained 2.3789, 10.5168, and
0.009528, while CA-RNN obtained 2.3474, 10.0759, and 0.007717. This
single-seed combined stress is exploratory evidence rather than an isolated
causal test of scale or dependence. Absolute levels and paired degradation are
considered together because a worse nominal starting point can mechanically
reduce measured deterioration.

\subsection{Financial results}

Table~\ref{tab:seed-results} reports averages over ten matched optimization seeds
and 957 forecast origins on the original financial test period. CA-RNN
reduced projected-PIT calibration error relative to RNN for all ten seeds, from
$0.001438$ to $0.001225$, or approximately 15\%. Its Energy Score and RMSE were
worse for all ten seeds. Its mean interval score was slightly lower, but the
paired difference remained unresolved, with a 95\% hierarchical interval of
$[-0.0191,0.0136]$. In contrast, its coverage-error difference was $-0.00528$
with interval $[-0.00815,-0.00136]$, favoring CA-RNN.

\begin{table}[ht]
\centering
\caption{Mean financial performance over ten matched optimization seeds. Lower
is better except for coverage, whose nominal target is 0.90.}
\label{tab:seed-results}
\resizebox{\textwidth}{!}{%
\begin{tabular}{lrrrrr}
\toprule
Model & RMSE & Energy & Interval & Coverage & PIT calibration \\
\midrule
RNN & 0.9399 & 1.1450 & 4.0855 & 0.8717 & 0.001438 \\
CA-RNN (temporal penalty removed) & 0.9504 & 1.1563 & 4.1191 & 0.8666 & 0.001306 \\
CA-RNN & 0.9479 & 1.1531 & 4.0816 & 0.8770 & 0.001225 \\
\bottomrule
\end{tabular}
}
\end{table}

The temporal-penalty-removal ablation isolates the contribution of that objective
component. Against RNN it had worse MSE,
Energy Score, interval score, and coverage error despite lower projected-PIT
error. CA-RNN then reduced MSE by $0.00465$, Energy Score by
$0.00317$, interval score by $0.03749$, and coverage error by $0.01032$ relative
to the temporal-penalty-removal ablation; every corresponding hierarchical
interval excluded zero. It repairs a
substantial part, but not all, of the cost of marginal calibration.

\subsection{Post-protocol financial evaluation audit}

The repeated original-panel fits were reevaluated with the exact projected
Gaussian PIT and held-out joint NLL. Table~\ref{tab:financial-audit} separates
these audit outcomes from the frozen sample-rank diagnostics. Analytic projected
calibration improves from $0.001439$ for RNN to $0.001223$ for CA--RNN, with
the latter lower for all ten optimization seeds. This agrees closely with the
sample-rank ordering ($0.001438$ versus $0.001225$), so the finite forecast-draw
PIT approximation does not explain the developmental calibration effect.
Nevertheless, CA--RNN has higher mean test NLL ($5.3095$ versus $5.2713$),
component-averaged MSE ($0.8986$ versus $0.8833$), and Energy Score ($1.1531$
versus $1.1450$). The variogram scores are nearly tied ($0.156705$ versus
$0.156821$), giving no clear dependence-sensitive proper-score advantage to
either method on the original panel. Analytic projected-PIT mean absolute
autocorrelation is higher for CA--RNN ($0.03691$ versus $0.03072$), despite the
temporal penalty. The audit thus strengthens, rather than reverses, the
calibration--predictive-quality distinction.

\begin{table}[ht]
\centering
\caption{Original-panel post-protocol evaluation audit: means across ten
matched neural seeds. All entries are losses or diagnostic discrepancies for
which lower is better. MSE and RMSE are distinct summaries.}
\label{tab:financial-audit}
\begin{tabular}{lrrrrr}
\toprule
Model & MSE & RMSE & Test NLL & Variogram & Analytic PIT error \\
\midrule
RNN & 0.8833 & 0.9399 & 5.2713 & 0.156821 & 0.001439 \\
CA--RNN & 0.8986 & 0.9479 & 5.3095 & 0.156705 & 0.001223 \\
\bottomrule
\end{tabular}
\end{table}

Using saved matched forecasts, a post-protocol block-length check repeated the
hierarchical seed-and-origin bootstrap with blocks of 10, 20, and 40 trading
days. The sign and zero-exclusion decision were stable across these lengths for
MSE (CA--RNN worse), Energy Score (worse), and sample-rank projected-PIT
calibration error (better). Interval-score and variogram-score differences
remained unresolved at all three lengths. Temporal PIT-dependence inference was
less reliable: the zero-exclusion decision changed between 10-day and longer
blocks, and the block-resampled interval was displaced from the observed point
difference. Artificial block boundaries can perturb lagged covariance, so this
bootstrap is not used as decisive evidence for temporal PIT adequacy. These
sensitivity analyses do not alter the frozen external-panel estimands or results.

\subsection{Conventional-baseline results}

Table~\ref{tab:baseline-results} adds the strongest traditional benchmark to the
original panel. Validation selected VAR lag one, and all four marginal GARCH
optimizations converged. The fixed-covariance Gaussian VAR was poorly calibrated;
adding marginal GARCH scales and resampling joint standardized residuals reduced
its interval score from 4.3993 to 4.0298 and projected-PIT error from 0.003388 to
0.000453. VAR--GARCH had better interval score, coverage error, and projected-PIT
calibration than CA-RNN. It also had slightly better RMSE and Energy
Score. The comparison attributes much of the neural models' remaining uncertainty
error to unmodeled conditional volatility rather than an insufficiently flexible
conditional mean alone.

\begin{table}[ht]
\centering
\caption{Original-panel neural and VAR--GARCH results. Neural entries are
ten-seed means; the traditional model is fitted once. Coverage targets 0.90.}
\label{tab:baseline-results}
\begin{tabular}{lrrrrr}
\toprule
Model & RMSE & Energy & Interval & Coverage & PIT calibration \\
\midrule
RNN & 0.9399 & 1.1450 & 4.0855 & 0.8717 & 0.001438 \\
CA-RNN & 0.9479 & 1.1531 & 4.0816 & 0.8770 & 0.001225 \\
VAR--GARCH bootstrap & 0.9452 & 1.1495 & \textbf{4.0298} & \textbf{0.8989} & \textbf{0.000453} \\
\bottomrule
\end{tabular}
\end{table}

Against CA-RNN, the paired hierarchical intervals favored
VAR--GARCH for MSE, Energy Score, interval score, coverage error, and projected-PIT
calibration. These intervals repeat the fixed traditional forecast across neural
seed slots; they represent test-origin dependence and neural optimization
variation but condition on the fitted baseline.

The training-sample specification checks in
Table~\ref{tab:var-garch-diagnostics} support the numerical and volatility
adequacy of the benchmark. The VAR companion spectral radius was 0.1513, well
below one. Every GARCH optimization converged, all persistence estimates were
below one, and no filtered variance reached the numerical floor. Persistence was
nevertheless high, ranging from 0.9700 to 0.9894, so volatility shocks decay
slowly. None of the 12 Ljung--Box checks on squared standardized residuals
rejected at the 5\% level, providing no evidence of remaining volatility
dependence at the tested lags.

\begin{table}[ht]
\centering
\caption{Training-sample diagnostics for the selected VAR(1)--GARCH benchmark.
$p_{z,20}$ and $p_{z^2,20}$ are descriptive Ljung--Box $p$-values at lag 20;
kurtosis is the ordinary, rather than excess, standardized-residual kurtosis.}
\label{tab:var-garch-diagnostics}
\begin{tabular}{lrrrr}
\toprule
Asset & $\alpha+\beta$ & $p_{z,20}$ & $p_{z^2,20}$ & Kurtosis \\
\midrule
AAPL & 0.9700 & 0.2307 & 0.9689 & 6.30 \\
JPM  & 0.9849 & 0.1490 & 0.7370 & 5.53 \\
XOM  & 0.9849 & 0.0142 & 0.8253 & 4.88 \\
WMT  & 0.9894 & 0.0089 & 0.8922 & 11.54 \\
\bottomrule
\end{tabular}
\end{table}

The conditional-mean diagnostics are less favorable. Six of the 12 Ljung--Box
tests on standardized residuals rejected at 5\%; XOM rejected at all three
tested lags, WMT at lags 10 and 20, and JPM at lag 5. Thus VAR(1) leaves some
linear serial dependence even though validation preferred it to the longer-lag
candidates. All four Jarque--Bera tests strongly rejected Gaussian standardized
innovations, and marginal kurtosis ranged from 4.88 to 11.54. Pairwise
standardized-residual correlations ranged from 0.22 to 0.49. The heavy tails and
cross-sectional dependence jointly support the use of synchronized empirical
residual vectors rather than Gaussian or independently resampled innovations.
Accordingly, VAR--GARCH is treated as a strong volatility-aware benchmark, not a
fully specified model of the conditional mean.

\subsection{Frozen external-panel replication}

The SHA-256 hash recorded with the executed results matches the frozen protocol.
The replication contains ten paired optimization
seeds and 989 test origins; all marginal GARCH optimizations converged.
Table~\ref{tab:confirmatory-results} reports model means.

\begin{table}[ht]
\centering
\caption{Frozen MSFT--BAC--CVX--COST replication. Neural entries are ten-seed
means; coverage targets 0.90.}
\label{tab:confirmatory-results}
\begin{tabular}{lrrrrrr}
\toprule
Model & RMSE & Energy & Interval & Coverage & PIT cal. & PIT dep. \\
\midrule
RNN & 0.8609 & 1.0677 & 3.7766 & 0.8703 & 0.000783 & 0.03008 \\
CA-RNN & 0.8694 & 1.0778 & 3.7998 & 0.8683 & 0.000735 & 0.03378 \\
VAR--GARCH bootstrap & 0.8646 & 1.0713 & 3.7528 & 0.8928 & 0.000491 & 0.03270 \\
\bottomrule
\end{tabular}
\end{table}

The frozen table retains the protocol's sample-rank projected PIT diagnostic so
that the prespecified primary outcome is not replaced after inspection. Exact
Gaussian projected PIT and test-NLL audits are reported for the development-panel
RNN and CA--RNN and for the replicated isolated-shift experiment, where the
conditional Gaussian parameters were retained. They are not reconstructed from
the external panel's saved forecast draws: samples alone do not identify the
exact predictive density. Likewise, no NLL is assigned to the empirical
VAR--GARCH residual bootstrap because it does not supply the same tractable
parametric density. This is a comparability boundary, not a zero or missing
performance result.

For the prespecified primary contrast, CA-RNN reduced projected-PIT
calibration error by $0.0000473$, with 95\% interval
$[-0.000163,-0.000012]$. This targeted effect therefore replicated on disjoint
assets. It did not yield a better predictive distribution overall:
candidate-minus-RNN differences were $+0.01476$ for MSE, $+0.01007$ for Energy
Score, and $+0.02318$ for interval score, with respective intervals
$[0.01236,0.01751]$, $[0.00847,0.01189]$, and $[0.00773,0.03937]$.
The coverage-error difference was unresolved. Projected PIT mean absolute
autocorrelation was also worse by $0.00370$, with interval
$[0.00005,0.00455]$. Thus the sequential objective did not improve out-of-sample
temporal PIT adequacy despite targeting finite-lag moments during training.

The secondary comparison again favored VAR--GARCH over CA-RNN for
MSE, Energy Score, interval score, coverage error, and projected calibration; all
five intervals excluded zero. The RNN retained better MSE and Energy Score than
VAR--GARCH, whereas VAR--GARCH had better interval score, coverage error, and
projected calibration. Conditional-mean accuracy and uncertainty calibration are
therefore supplied by different model components in this application.

\subsection{Rolling market-regime results}

Table~\ref{tab:regime-results} gives five-seed means. CA-RNN had lower
projected-PIT calibration error than RNN in all four regimes. Hierarchical
intervals favored it in the pre-pandemic, inflation/tightening, and recent
windows; the pandemic difference was unresolved. In contrast, its Energy Score
was lower only pre-pandemic, where the difference was unresolved.

\begin{table}[ht]
\centering
\caption{Five-seed regime means for RNN (R) and CA-RNN (C).}
\label{tab:regime-results}
\begin{tabular}{lrrrrrr}
\toprule
& \multicolumn{2}{c}{Energy Score} & \multicolumn{2}{c}{Interval score}
& \multicolumn{2}{c}{PIT calibration} \\
Regime & R & C & R & C & R & C \\
\midrule
Pre-pandemic & 0.8629 & 0.8626 & 3.1123 & 3.1027 & 0.001914 & 0.001503 \\
Pandemic & 1.5016 & 1.5029 & 5.7371 & 5.6629 & 0.002995 & 0.002883 \\
Inflation/tightening & 1.2809 & 1.2884 & 4.5862 & 4.5396 & 0.001842 & 0.001465 \\
Recent & 1.0237 & 1.0289 & 3.6648 & 3.6935 & 0.001224 & 0.000904 \\
\bottomrule
\end{tabular}
\end{table}

The interval-score ordering is regime-dependent. CA-RNN significantly
outperformed RNN during the pandemic and inflation/tightening periods, was
unresolved pre-pandemic, and significantly underperformed recently. Coverage
error similarly favored CA-RNN pre-pandemic and during
inflation/tightening, was unresolved during the pandemic, and favored RNN
recently. These reversals directly answer the distribution-shift question.

Across regimes, CA-RNN improved Energy Score and interval score over the
temporal-penalty-removal ablation directionally in four of four windows. These improvements were
resolved in the pandemic, inflation/tightening, and recent windows, while the
pre-pandemic differences were unresolved. Projected calibration improved over
the temporal-penalty-removal ablation in three of four windows.

\begin{figure}[ht]
\centering
\IfFileExists{../results/market_regimes/regime_metric_profiles.png}{%
\includegraphics[width=\textwidth]{../results/market_regimes/regime_metric_profiles.png}%
}{\fbox{\parbox{0.9\textwidth}{Run notebook 09 to generate the rolling-regime
metric profile.}}}
\caption{Mean Energy Score, interval score, and projected-PIT calibration error
across five matched seeds and four chronological test regimes. Lower is better.}
\label{fig:regime-profiles}
\end{figure}

\subsection{Interval-score mechanism under temporal shift}

To explain the regime reversal, the retained forecasts were analyzed without
refitting. For a 90\% interval $[L_{t,j},U_{t,j}]$, the componentwise interval
score decomposes exactly as
\begin{equation}
 S_{t,j}=\underbrace{U_{t,j}-L_{t,j}}_{\text{width}}
 +\underbrace{20(L_{t,j}-Y_{t,j})\ind\{Y_{t,j}<L_{t,j}\}}_{\text{lower-tail penalty}}
 +\underbrace{20(Y_{t,j}-U_{t,j})\ind\{Y_{t,j}>U_{t,j}\}}_{\text{upper-tail penalty}}.
 \label{eq:interval-decomposition}
\end{equation}
Table~\ref{tab:interval-mechanism} reports the mean CA-RNN-minus-RNN
difference in each component. The arithmetic identity in
\eqref{eq:interval-decomposition} means the final column is exactly the sum of the
preceding three columns.

\begin{table}[ht]
\centering
\caption{Decomposition of the interval-score difference between CA-RNN and RNN.
Negative total differences favor CA-RNN.}
\label{tab:interval-mechanism}
\begin{tabular}{lrrrr}
\toprule
Regime & Width & Lower penalty & Upper penalty & Total \\
\midrule
Pre-pandemic & $-0.0439$ & $+0.0241$ & $+0.0103$ & $-0.0096$ \\
Pandemic & $+0.0420$ & $-0.0599$ & $-0.0563$ & $-0.0742$ \\
Inflation/tightening & $+0.2342$ & $-0.1709$ & $-0.1098$ & $-0.0466$ \\
Recent & $-0.0185$ & $+0.0268$ & $+0.0205$ & $+0.0287$ \\
\bottomrule
\end{tabular}
\end{table}

The pandemic and inflation/tightening improvements arise through beneficial
widening: the added width is more than offset by reductions in both tail-miss
penalties. In the recent period, CA-RNN instead produces slightly
narrower intervals and larger penalties in both tails. Before the pandemic, the
mechanism is different again: reduced width outweighs the additional tail
penalties, yielding a small and statistically unresolved sharpness benefit.

The stress-period effect is concentrated in forecasts preceded by elevated
volatility. CA-RNN's interval-score differences in the pandemic are
$+0.0198$, $-0.0725$, and $-0.1705$ for low-, medium-, and high-volatility
origins, respectively. During inflation/tightening they are $+0.0209$, $-0.0562$,
and $-0.1045$. By contrast, the recent-period differences are positive in every
state: $+0.0318$, $+0.0238$, and $+0.0305$. Volatility states use trailing 20-day
realized volatility known before each forecast and within-regime terciles.

The improvements are also asset-concentrated. XOM accounts for most of the
inflation/tightening gain, with an interval-score difference of $-0.2364$;
AAPL, JPM, and WMT have differences of $-0.0051$, $+0.0221$, and $+0.0331$.
Pandemic improvements occur for AAPL, XOM, and WMT but not JPM. Recent losses are
concentrated in AAPL, XOM, and WMT, while JPM improves slightly. These results
preclude interpreting an aggregate gain as uniform across assets.

The train-to-test descriptors support, but do not prove, a scale-and-tail
interpretation. Pandemic volatility is 1.26 times its training level and the mean
standardized fourth moment is 29.87. The corresponding values are 1.00 and 12.08
during inflation/tightening, versus 0.83 and 6.37 recently. The recent period has
the largest standardized location shift and correlation shift of the four
windows. Thus the complete CA-RNN objective appears most useful for persistent
scale and tail errors, while lower-volatility location or dependence changes
remain harder.

Importantly, projected and coordinate PIT calibration still improve recently even
as 90\% interval performance deteriorates. The calibration loss averages 19 CDF
grid levels and multiple projections; it does not target the 5\% and 95\% quantiles
specifically. Better global PIT uniformity can therefore coexist with worse tail
coverage. The exact score decomposition is established by the retained forecasts,
whereas the proposed economic mechanism is post-hoc and requires new data for
confirmation.

\begin{figure}[ht]
\centering
\IfFileExists{../results/interval_mechanisms/interval_score_decomposition.png}{%
\includegraphics[width=0.82\textwidth]{../results/interval_mechanisms/interval_score_decomposition.png}%
}{\fbox{\parbox{0.8\textwidth}{Run notebook 10 to generate the interval-score
decomposition.}}}
\caption{Exact decomposition of CA-RNN-minus-RNN interval-score
differences. Tail reductions offset widening in the pandemic and
inflation/tightening periods, but not in the recent period.}
\label{fig:interval-mechanism}
\end{figure}

\section{Discussion}

\subsection{Interpretation}

Four conclusions follow. First, a marginal calibration penalty is insufficient
for dependent multivariate forecasting: it can improve its targeted diagnostic
while harming proper scores and nominal coverage. Second, the temporal PIT-moment
loss produces a genuine but narrow intervention on serial PIT moments. It
recovered part of the temporal-penalty-removal model's predictive-performance
loss in the ten-seed financial experiment, but the refreshed controlled
simulation found only a small temporal-dependence reduction accompanied by worse
point error, proper scores, projected calibration, and coverage error. Its value
is therefore data- and regime-dependent rather than automatic. Third, CA-RNN's projected-calibration
effect transfers to a disjoint asset panel, but its developmental interval and
coverage improvements do not: the frozen replication instead favors RNN on all
three proper or accuracy-oriented scores and finds no resolved coverage benefit.
Fourth, VAR--GARCH shows that explicit volatility dynamics and joint residual
resampling can provide better-calibrated uncertainty than direct calibration
regularization in these data.

The replicated isolated-scale experiment provides a complementary and more
specific robustness finding. Across 50 independent Gaussian DGP realizations,
CA--RNN offers little nominal proper-score advantage but degrades less than RNN
when only test innovation scale rises. Its shifted gains span joint NLL, Energy
Score, variogram score, marginal interval score, and analytic projected-PIT
calibration, rather than a single diagnostic. This does not contradict the
frozen financial replication: the simulation isolates one known perturbation
of a Gaussian autoregression, while the external asset panel combines unknown
volatility, tail, dependence, and mean changes. The post-protocol simulation
therefore establishes a boundary at which the fixed CA--RNN objective helps,
not a general claim of financial forecast superiority.

The disagreement between Energy Score and interval score is substantively useful.
Energy Score assesses the multivariate predictive distribution, whereas the
marginal interval score rewards coverage and sharpness asset by asset. During the
pandemic and inflation/tightening periods, CA-RNN improves the latter
without improving the former. Its practical value therefore depends on whether an
application prioritizes marginal risk intervals or overall multivariate accuracy.
The confirmatory evidence narrows this claim further: lower projected-PIT grid
error alone is insufficient evidence of improved operational uncertainty.

The regret identities explain the observed divergence. Energy Score regret is
one half of the energy distance and therefore responds to discrepancies anywhere
in the joint distribution. Marginal interval-score regret is the sum of two
integrated quantile errors at the selected tail probabilities. By contrast, the
training regularizer averages a finite set of smoothed PIT-CDF moments across
grid points and projections. Because this finite criterion is non-identifying,
reducing it supplies no mathematical upper bound on either regret. The external
result---lower projected-PIT error alongside higher Energy and interval
scores---is therefore compatible with the objectives rather than paradoxical.

\subsection{Limitations}

The central limitation concerns what the calibration objective can identify.
CA-RNN evaluates a finite set of smoothed projected-PIT moments, estimated from
dependent mini-batches. Its conclusions therefore depend on the selected
projections, grid, smoothing parameter, and batching scheme. Even a small value
of this empirical loss cannot establish full joint or history-conditional
calibration. The temporal term narrows one part of this gap by targeting selected
lagged covariances, but finite zero autocovariances do not imply PIT independence.
Moving-block resampling introduces an additional complication because artificial
block boundaries can affect autocorrelation estimates. The 10/20/40-day
sensitivity audit changed the zero-exclusion decision for PIT dependence and
showed resampling bias relative to its observed point difference; the
corresponding uncertainty intervals are therefore interpreted as diagnostic
rather than definitive.

The empirical scope is also deliberately limited. The financial evidence comes
from two four-asset panels and one-step-ahead forecasts, while the chronological
regimes are calendar-based and were partly formulated after broader exposure to
the development data. The frozen external panel reduces, but does not eliminate,
dependence on choices of assets, architecture, tuning path, projections, and
evaluation windows. Similarly, the hierarchical bootstrap represents variation
across neural optimization seeds and dependent forecast origins conditional on
those design choices. Because each conventional benchmark is fitted once per
panel, its reported comparisons do not integrate over baseline parameter-
estimation uncertainty. The isolated-scale simulation adds 50 independent
datasets, but uses one neural seed per dataset and was designed after earlier
developmental findings. Its MCSE therefore does not separate DGP from
optimization variation or protect against post-selection of the shift family.

Finally, the neural predictive law is Gaussian and does not explicitly encode
the persistent conditional volatility or heavy tails characteristic of financial
returns. This matters for interpretation: VAR--GARCH's advantage may reflect a
better uncertainty specification rather than a general advantage of linear over
recurrent conditional-mean dynamics. The original-panel benchmark is itself not
fully dynamically specified. Its diagnostics are in-sample and descriptive,
several Ljung--Box tests on standardized residual levels reject, and the
$p$-values are not adjusted for fitted parameters or multiple testing. Its strong
performance is therefore attributed primarily to effective volatility filtering
and joint empirical innovation resampling, not to complete model adequacy.

\subsection{Future work}

The most direct methodological extension is a tail-aware sequential calibration
objective. The current loss weights all 19 PIT grid levels and all fixed
projections equally. A weighted version could emphasize application-relevant tail
levels,
\begin{equation}
 \mathcal L_{\mathrm{tail}}
 =\frac{1}{R\sum_g w_g}\sum_{r=1}^R\sum_{g=1}^G w_g
 \left[\frac{1}{|B|}\sum_{t\in B}s_\tau(q_g-U_{t,r})-q_g\right]^2,
 \label{eq:tail-calibration}
\end{equation}
with larger prespecified $w_g$ near $q=0.05$ and $q=0.95$. A complementary
extension would penalize lagged dependence in tail exceedance indicators rather
than only covariance of centered PITs. Either extension must be tuned on
validation data and evaluated with global calibration and proper scores to avoid
improving one nominal interval at the expense of the rest of the distribution.

A particularly relevant extension is a volatility-aware CA-RNN. The recurrent
state could parameterize time-varying marginal scales and a structured dynamic
correlation matrix, with Student-$t$ or mixture innovations replacing the current
Gaussian head. This would test whether calibration-aware estimation adds value
after conditional heteroskedasticity and heavy tails have been modeled directly,
rather than requiring the calibration penalty to compensate for their omission.
It is intentionally left outside the present study because a defensible
evaluation would require a factorial comparison of volatility-aware RNN and
volatility-aware CA-RNN, alongside their current Gaussian counterparts. Without
that control, improvements could not be attributed separately to volatility
modeling, the innovation distribution, or calibration-aware training.

Further work should also test new frequencies, forecast horizons, broader asset
universes, and genuinely future calendar periods. Conditional calibration
diagnostics could use volatility, market return, and liquidity states fixed
before evaluation. Finally, uncertainty for PIT dependence should use a
procedure that avoids correlations introduced at moving-block boundaries.

\section{Conclusion}

This paper adapts calibration-aware Bayesian learning to continuous multivariate
recurrent forecasting and introduces a temporal calibration penalty based on
lagged projected-PIT covariance. Calibration-aware training reliably changes
forecast calibration, but projected calibration alone incurs predictive costs.
The temporal PIT-moment component improves the corresponding loss-component
ablation in the ten-seed financial study, but not across every metric in the
controlled simulation. The complete CA-RNN nevertheless produces lower
projected-PIT calibration error than a likelihood-trained RNN across financial
optimization seeds, chronological regimes, and a frozen disjoint-asset
replication; the same ordering does not hold in the single-realization nonlinear
simulation.

The method is not uniformly superior. In the confirmatory experiment it has worse
MSE, Energy Score, interval score, and temporal PIT dependence than RNN, while its
coverage-error difference is unresolved. VAR--GARCH also outperforms it on the
principal uncertainty metrics in both asset panels. The contribution is therefore
not a dominance claim, nor evidence that the temporal PIT-moment penalty achieves
PIT independence. It is a controlled demonstration that calibration-aware
training can reproducibly improve its targeted finite-projection diagnostic while
failing to improve---and sometimes worsening---proper scores and sequential
adequacy. In the completed post-protocol isolated-scale experiment, 50 paired
realizations show a more precise robustness crossover: nominal near-parity gives
way to lower NLL, Energy Score, variogram score, interval score, and analytic
projected-PIT error for CA--RNN as test innovation scale rises. This is evidence
for one isolated stress mechanism, not general dominance under distribution
shift. Strong volatility-aware
baselines, ablation across data-generating settings, and external confirmation
are thus essential when evaluating calibration-aware neural forecasts.

\newpage
\appendix

\section{Reproducibility and post-protocol evaluation audit}

\subsection{Data, transformations, and chronological partitions}

The development panel contains adjusted daily closing prices for AAPL, JPM,
XOM, and WMT, requested for 2007-01-01 through 2026-01-01 (exclusive). After
complete-case date alignment and one-day log differencing, it contains 4,779
return observations from 2007-01-04 through 2025-12-31. The 60/20/20
chronological split gives training (2,867 observations, 2007-01-04 to
2018-05-23), validation (955, 2018-05-24 to 2022-03-09), and test (957,
2022-03-10 to 2025-12-31). The observed adjusted-price span begins
2007-01-03. The cleaning rule rejects any ticker with over 1\% missingness,
then retains complete dates and positive prices. The panel is subject to
survivorship-biased asset selection.

The frozen external panel contains MSFT, BAC, CVX, and COST in that order.
Its 4,945 aligned log-return observations run from 2007-01-04 through
2026-08-31. The same split gives training (2,967 observations, 2007-01-04 to
2018-10-15), validation (989, 2018-10-16 to 2022-09-20), and test (989,
2022-09-21 to 2026-08-31). Each component is standardized using its own
training-partition mean and sample standard deviation. A 20-observation
historical context precedes each validation and test target; neither scaler
nor fitted model sees the target or later observations during selection.

\subsection{Model, randomization, and selection settings}

The focused neural architecture is a one-layer, 32-hidden-unit GRU with a
full-covariance Gaussian predictive head and minimum Cholesky diagonal scale
$10^{-4}$. It uses a 20-day input history, AdamW learning rate $10^{-3}$,
batch size 64, maximum 200 epochs, patience 25, and validation-NLL checkpoint
selection in the repeated financial experiments. The CA--RNN loss weights are
$\lambda_{\mathrm{cal}}=100$ and $\lambda_{\mathrm{seq}}=2000$; the PIT CDF grid has
19 points, smoothing temperature 0.03, and five temporal lags. Projection
directions comprise four coordinates, an equal-weight direction, and four
random unit vectors generated with seed 606. The repeated financial and frozen
external studies use neural training seeds 611, 619, 631, 641, 647, 653, 661,
673, 683, and 691, forecast seed 20,606, and 500 draws per forecast origin.
The original-panel hierarchical bootstrap uses 20-trading-day blocks, 2,000
draws, and seed 808; the frozen panel uses the same block length and draw count
with seed 1208. All paired comparisons reuse forecast-origin indices across
models. The conventional benchmark selects VAR lag from $\{1,5,10,20\}$ by
validation one-step MSE and then fits marginal GARCH$(1,1)$ and the joint
standardized-residual resampling law on the training partition.

The four rolling test windows are 2017--2019, 2020--2021, 2022--2023, and
2024--2025. Their expanding training endpoints are, respectively, the ends of
2014, 2016, 2019, and 2021; intervening validation periods are 2015--2016,
2017--2019, 2020--2021, and 2022--2023. The regime training seeds are
701, 709, 719, 727, and 733. Their
1,000-draw hierarchical bootstraps are developmental rather than frozen
confirmatory analyses. The four-model nonlinear demonstration instead uses
1,200 observations, DGP and training seed 2027, 150 maximum epochs, and patience
20. The original mixed scale-and-correlation curve uses DGP and training seed
505, 120 maximum epochs, and patience 15; its single-realization interpretation
remains exploratory.

The completed post-protocol isolated-scale experiment is implemented in notebook
14 and reported in Table~\ref{tab:isolated-scale}. It uses 50 independent
Gaussian autoregressive realizations with DGP seeds 1000--1049, one distinct
neural seed $606+b$ per realization $b=0,\ldots,49$, and common forecast seed
20,606. Each series has 1,200 observations; the first 60\% train, the next
20\% validate, and the final 20\% test. The shift begins exactly at the test
boundary. Severities $\{0,0.5,1\}$ multiply the post-boundary innovation
standard deviation by $1+\delta$ while holding correlation, autoregressive
coefficients, and innovation family fixed. RNN and CA--RNN are fitted once per
realization with 120 maximum epochs, patience 15, and validation-NLL
checkpointing on the same unshifted training/validation history, then held fixed
over severities. The run records failures and training time. For each severity,
paired CA--RNN-minus-RNN differences are averaged across valid realizations,
with Monte Carlo standard error $s_D/\sqrt{B}$. Prefixes of 10, 25, and 50
realizations provide a replication-count sensitivity check. One neural seed
per DGP realization means DGP and optimization variability are not separately
identified in this experiment.

\subsection{Evaluation definitions and protocol status}

For a $d$-asset standardized target $y_t$ and forecast mean $\widehat\mu_t$,
the reported component-averaged mean squared error is
$\mathrm{MSE}=(Td)^{-1}\sum_{t,j}(y_{t,j}-\widehat\mu_{t,j})^2$ and
$\mathrm{RMSE}=\sqrt{\mathrm{MSE}}$. A seed-level mean of RMSEs need not equal
the square root of the corresponding seed-level mean MSE. Neural test NLL is
the mean joint Gaussian negative log density, evaluated at held-out targets;
for Bayesian models it is the negative log of a Monte Carlo predictive mixture.
For a deterministic Gaussian forecast, projected PITs are evaluated analytically
as $\Phi[(a^\top y_t-a^\top\mu_t)/\sqrt{a^\top\Sigma_t a}]$ in the new
evaluation audit. The previously reported sample-rank PITs remain the frozen
protocol measure and are not retrospectively replaced. The variogram score is
reported as the mean over origins and unordered asset pairs of
$(|y_{t,i}-y_{t,j}|^{1/2}-E_F|Y_i-Y_j|^{1/2})^2$; the expectation is estimated
from forecast draws. Notebook 15 checks sensitivity of developmental paired
intervals to block lengths 10, 20, and 40, holding the forecast arrays fixed.
These additions are secondary, post-protocol audits, not new primary outcomes.

The frozen primary protocol is preserved verbatim at
\texttt{docs/confirmatory\_protocol.md}; its SHA-256 digest is
\nolinkurl{20a0531461992b11d702d5b72b51c57a5e9b162e84dd5f6e6e34e5fadec570b0}.
Its historical ``sequential CA--RNN'' label denotes the unchanged full
projected-plus-temporal objective now called CA--RNN. The protocol fixes the
external assets and date boundary, architecture, penalties, seeds, forecast
draws, primary metrics, and 20-day hierarchical bootstrap. No new analytic
PIT, NLL, variogram, or sensitivity result supersedes its originally frozen
sample-PIT and proper-score comparisons.

\begin{thebibliography}{99}

\bibitem{huang2024cabl}
J. Huang, S. Park, and O. Simeone,
``Calibration-Aware Bayesian Learning,''
\emph{arXiv preprint arXiv:2305.07504}, version 2, 2024.
\href{https://arxiv.org/abs/2305.07504}{arXiv:2305.07504}.

\bibitem{gneiting2007calibration}
T. Gneiting, F. Balabdaoui, and A. E. Raftery,
``Probabilistic Forecasts, Calibration and Sharpness,''
\emph{Journal of the Royal Statistical Society: Series B},
vol. 69, no. 2, pp. 243--268, 2007.
\href{https://doi.org/10.1111/j.1467-9868.2007.00587.x}{doi:10.1111/j.1467-9868.2007.00587.x}.

\bibitem{gneiting2007scoring}
T. Gneiting and A. E. Raftery,
``Strictly Proper Scoring Rules, Prediction, and Estimation,''
\emph{Journal of the American Statistical Association},
vol. 102, no. 477, pp. 359--378, 2007.
\href{https://doi.org/10.1198/016214506000001437}{doi:10.1198/016214506000001437}.

\bibitem{gneiting2008multivariate}
T. Gneiting, L. I. Stanberry, E. P. Grimit, L. Held, and N. A. Johnson,
``Assessing Probabilistic Forecasts of Multivariate Quantities, with an
Application to Ensemble Predictions of Surface Winds,''
\emph{TEST}, vol. 17, pp. 211--235, 2008.
\href{https://doi.org/10.1007/s11749-008-0114-x}{doi:10.1007/s11749-008-0114-x}.

\bibitem{scheuerer2015variogram}
M. Scheuerer and T. M. Hamill,
``Variogram-Based Proper Scoring Rules for Probabilistic Forecasts of
Multivariate Quantities,''
\emph{Monthly Weather Review}, vol. 143, no. 4, pp. 1321--1334, 2015.
\href{https://doi.org/10.1175/MWR-D-14-00269.1}{doi:10.1175/MWR-D-14-00269.1}.

\bibitem{gneiting2023conditional}
T. Gneiting and J. Resin,
``Regression Diagnostics Meets Forecast Evaluation: Conditional Calibration,
Reliability Diagrams, and Coefficient of Determination,''
\emph{Electronic Journal of Statistics}, vol. 17, no. 2,
pp. 3226--3286, 2023.
\href{https://doi.org/10.1214/23-EJS2180}{doi:10.1214/23-EJS2180}.

\bibitem{diebold1998density}
F. X. Diebold, T. A. Gunther, and A. S. Tay,
``Evaluating Density Forecasts, with Applications to Financial Risk
Management,'' \emph{International Economic Review}, vol. 39, no. 4,
pp. 863--883, 1998.
\href{https://www.nber.org/papers/t0215}{NBER working-paper record}.

\bibitem{berkowitz2001testing}
J. Berkowitz,
``Testing Density Forecasts, with Applications to Risk Management,''
\emph{Journal of Business \& Economic Statistics}, vol. 19, no. 4,
pp. 465--474, 2001.
\href{https://doi.org/10.1198/07350010152596718}{doi:10.1198/07350010152596718}.

\bibitem{knuppel2015multistep}
M. Kn\"uppel,
``Evaluating the Calibration of Multi-Step-Ahead Density Forecasts Using Raw
Moments,'' \emph{Journal of Business \& Economic Statistics}, vol. 33,
no. 2, pp. 270--281, 2015.
\href{https://ideas.repec.org/a/taf/jnlbes/v33y2015i2p270-281.html}{publication record}.

\bibitem{guo2017calibration}
C. Guo, G. Pleiss, Y. Sun, and K. Q. Weinberger,
``On Calibration of Modern Neural Networks,'' in
\emph{Proceedings of the 34th International Conference on Machine Learning},
vol. 70, pp. 1321--1330, 2017.
\href{https://proceedings.mlr.press/v70/guo17a.html}{PMLR article page}.

\bibitem{kumar2018trainable}
A. Kumar, S. Sarawagi, and U. Jain,
``Trainable Calibration Measures for Neural Networks from Kernel Mean
Embeddings,'' in \emph{Proceedings of the 35th International Conference on
Machine Learning}, vol. 80, pp. 2805--2814, 2018.
\href{https://proceedings.mlr.press/v80/kumar18a.html}{PMLR article page}.

\bibitem{kuleshov2018regression}
V. Kuleshov, N. Fenner, and S. Ermon,
``Accurate Uncertainties for Deep Learning Using Calibrated Regression,'' in
\emph{Proceedings of the 35th International Conference on Machine Learning},
vol. 80, pp. 2796--2804, 2018.
\href{https://proceedings.mlr.press/v80/kuleshov18a.html}{PMLR article page}.

\bibitem{utpala2020quantile}
S. Utpala and P. Rai,
``Quantile Regularization: Towards Implicit Calibration of Regression Models,''
\emph{arXiv preprint arXiv:2002.12860}, 2020.
\href{https://arxiv.org/abs/2002.12860}{arXiv:2002.12860}.

\bibitem{dheur2023large}
J. Dheur and S. Ben Taieb,
``A Large-Scale Study of Probabilistic Calibration in Neural Network Regression,''
in \emph{Proceedings of the 40th International Conference on Machine Learning},
vol. 202, pp. 7813--7836, 2023.
\href{https://proceedings.mlr.press/v202/dheur23a.html}{PMLR article page}.

\bibitem{dheur2024design}
J. Dheur and S. Ben Taieb,
``Probabilistic Calibration by Design for Neural Network Regression,''
in \emph{Proceedings of the 27th International Conference on Artificial
Intelligence and Statistics}, vol. 238, pp. 3133--3141, 2024.
\href{https://proceedings.mlr.press/v238/dheur24a.html}{PMLR article page}.

\bibitem{allen2024multivariate}
S. Allen, J. F. Ziegel, and D. Ginsbourger,
``Assessing the Calibration of Multivariate Probabilistic Forecasts,''
\emph{Quarterly Journal of the Royal Meteorological Society}, vol. 150,
no. 760, pp. 1315--1335, 2024.
\href{https://doi.org/10.1002/qj.4647}{doi:10.1002/qj.4647}.

\bibitem{chung2021pinball}
Y. Chung, W. Neiswanger, I. Char, and J. Schneider,
``Beyond Pinball Loss: Quantile Methods for Calibrated Uncertainty
Quantification,'' in \emph{Advances in Neural Information Processing Systems},
vol. 34, 2021.
\href{https://proceedings.neurips.cc/paper/2021/hash/5b168fdba5ee5ea262cc2d4c0b457697-Abstract.html}{NeurIPS article page}.

\bibitem{desobry2025prerank}
N. Desobry, E. Zhalieva, and S. Ben Taieb,
``Enforcing Calibration in Multi-Output Probabilistic Regression with Pre-rank
Regularization,'' \emph{arXiv preprint arXiv:2510.21273}, 2025.
\href{https://arxiv.org/abs/2510.21273}{arXiv:2510.21273}.

\bibitem{laajil2026prerank}
A. Laajil, E. Zhalieva, N. Desobry, and S. Ben Taieb,
``Calibrated Multivariate Distributional Regression with Pre-Rank
Regularization,'' \emph{arXiv preprint arXiv:2601.22895}, 2026.
\href{https://arxiv.org/abs/2601.22895}{arXiv:2601.22895}.

\bibitem{salinas2020deepar}
D. Salinas, V. Flunkert, J. Gasthaus, and T. Januschowski,
``DeepAR: Probabilistic Forecasting with Autoregressive Recurrent Networks,''
\emph{International Journal of Forecasting}, vol. 36, no. 3,
pp. 1181--1191, 2020.
\href{https://doi.org/10.1016/j.ijforecast.2019.07.001}{doi:10.1016/j.ijforecast.2019.07.001}.

\bibitem{white1982misspecification}
H. White,
``Maximum Likelihood Estimation of Misspecified Models,''
\emph{Econometrica}, vol. 50, no. 1, pp. 1--25, 1982.
\href{https://doi.org/10.2307/1912526}{doi:10.2307/1912526}.

\bibitem{ovadia2019shift}
Y. Ovadia et al.,
``Can You Trust Your Model's Uncertainty? Evaluating Predictive Uncertainty
under Dataset Shift,'' in \emph{Advances in Neural Information Processing
Systems}, vol. 32, 2019.
\href{https://proceedings.neurips.cc/paper_files/paper/2019/hash/8558cb408c1d76621371888657d2eb1d-Abstract.html}{NeurIPS article page}.

\bibitem{bollerslev1986garch}
T. Bollerslev,
``Generalized Autoregressive Conditional Heteroskedasticity,''
\emph{Journal of Econometrics}, vol. 31, no. 3, pp. 307--327, 1986.
\href{https://doi.org/10.1016/0304-4076(86)90063-1}{doi:10.1016/0304-4076(86)90063-1}.

\end{thebibliography}

\end{document}
